\documentclass[11pt,a4paper]{report}
\usepackage[margin=1in]{geometry}
\usepackage{amsmath,amssymb,mathtools,bm}
\usepackage{booktabs}
\usepackage{array}
\usepackage{longtable}
\usepackage{enumitem}
\usepackage{hyperref}
\usepackage{xcolor}
\usepackage{listings}
\usepackage{microtype}
\usepackage[T1]{fontenc}
\usepackage[utf8]{inputenc}

\hypersetup{colorlinks=true,linkcolor=blue,citecolor=blue,urlcolor=blue}
\lstset{basicstyle=\ttfamily\small,breaklines=true,columns=fullflexible,frame=single}

\newcommand{\dd}{\mathrm{d}}
\newcommand{\ii}{\mathrm{i}}
\newcommand{\eps}{\varepsilon}
\newcommand{\LTD}{\operatorname{LTD}}
\newcommand{\CFF}{\operatorname{CFF}}
\newcommand{\OSE}{\operatorname{OSE}}
\newcommand{\Res}{\operatorname{Res}}
\newcommand{\Graph}{\mathcal G}
\newcommand{\Basis}{\mathcal B}
\newcommand{\Orient}{\mathcal O}
\newcommand{\Surf}{\mathcal S}
\newcommand{\Weights}{\mathcal W}
\newcommand{\Active}{\mathcal A}
\newcommand{\Samples}{\mathcal X}
\newcommand{\Ints}{\mathcal E}
\newcommand{\Rep}{\mathcal R}
\newcommand{\Pinch}{\mathcal P}
\newcommand{\rank}{\operatorname{rank}}

\title{Generalised Three-Dimensional LTD/CFF Representations}
\author{Implementation and theory note}
\date{April 2026}

\begin{document}
\maketitle

\begin{abstract}
This note describes the derivative-free three-dimensional representation
implemented in \texttt{generalised\_ltd.py}.  The construction combines loop-tree
duality (LTD) with the causal/CFF representation in order to handle repeated
propagators while keeping the numerator as a sampled black box.  It also
describes the bounded-energy-degree extension used when edge-momentum
representation (EMR) numerators contain powers above the affine CFF class.

The document is organized as follows.  Chapter~1 gives the mathematical
construction, with all dependencies made explicit.  Chapter~2 explains how the
Python implementation and the oriented JSON schema encode the formulae.
Chapter~3 reports current Symbolica evaluator timings from the example scripts.
\end{abstract}

\tableofcontents

\chapter{Theoretical Construction}

\section{Framing, Notation, LTD, and CFF}

Let \(\Graph\) be a connected scalar graph with \(L\) loop momenta and internal
edge set \(E\).  A loop-momentum basis (LMB) is denoted by
\(\Basis=(k_1,\ldots,k_L)\).  For an edge \(e\in E\),
\begin{equation}
  q_e(k,p)=\sum_{r=1}^{L}\rho_{er}k_r+\sum_{a=1}^{n_{\rm ext}}\eta_{ea}p_a,
  \qquad \rho_{er},\eta_{ea}\in\mathbb Z ,
\end{equation}
and
\begin{equation}
  q_e^0(k^0,p^0)=\rho_e\cdot k^0+\eta_e\cdot p^0,\qquad
  \omega_e(\vec k,p,m)=\sqrt{\vec q_e^{\,2}+m_e^2}.
\end{equation}
The Feynman denominator is
\begin{equation}
  D_e=(q_e^0)^2-\omega_e^2+\ii0 .
\end{equation}

For a numerator \(N\) in EMR variables, the \(D\)-dimensional input integral is
\begin{equation}
  I_{\Graph,\Basis}[N]
  =
  \int_{\mathbb R^{4L}}
  \prod_{r=1}^{L}\frac{\dd^4 k_r}{(2\pi)^4}\,
  N(\{q_e(k,p)\}_{e\in E})
  \prod_{e\in E}D_e^{-a_e}.
  \label{eq:four-dimensional-integral}
\end{equation}
The powers \(a_e\) are normally one; repeated propagators are encoded below as
explicit repeated edges or, equivalently, as collapsed channels with powers.

\paragraph{LTD dependency data.}
The LTD construction is an algorithmic residue expansion of
\eqref{eq:four-dimensional-integral} after choosing
\begin{itemize}[leftmargin=2em]
  \item the graph \(\Graph\),
  \item the loop-momentum basis \(\Basis\),
  \item a contour-closure prescription \(\kappa\), i.e. the side on which each
        loop-energy integration is closed.
\end{itemize}
We denote the resulting spatial integrand by
\begin{equation}
  I^{(3)}_{\LTD}[\Graph,\Basis,\kappa;N](\vec k;p,m)
  =
  \sum_{\mathfrak o\in\Orient_{\LTD}(\Graph,\Basis,\kappa)}
  N\!\left(Q_{\mathfrak o}^{\LTD}(\vec k;p,m)\right)
  W_{\mathfrak o}^{\LTD}(\vec k;p,m).
\end{equation}
The map \(Q_{\mathfrak o}^{\LTD}\) gives every EMR edge energy after solving
the cut equations.  The weight \(W_{\mathfrak o}^{\LTD}\) is a product of
cut half-edge factors and LTD dual surfaces.  In this note we do not rederive
LTD; we use the multiloop algorithmic construction of
Ref.~\cite{ltd1906}.

\paragraph{CFF dependency data.}
For CFF/causal representations we use the graph as the primary input and assume
the canonical source/sink contraction ordering used by the implementation.
Thus
\begin{equation}
  I^{(3)}_{\CFF}[\Graph;N](\vec k;p,m)
  =
  \sum_{\mathfrak o\in\Orient_{\CFF}(\Graph)}
  N\!\left(Q_{\mathfrak o}^{\CFF}(\vec k;p,m)\right)
  W_{\mathfrak o}^{\CFF}(\vec k;p,m).
\end{equation}
The CFF weights use causal \(E\)-surfaces in the denominator.  The construction
is the one of Refs.~\cite{cff2211,cff2311}.  The standard black-box numerator
assumption is that \(N\) is at most affine in each EMR edge energy \(q_e^0\).

\paragraph{Surface notation.}
For any integer vector \(c_e\in\{-1,0,1\}\) and external-energy shift
\(\Delta\), define the signed causal surface
\begin{equation}
  \eta(c,\Delta)=\sum_{e}c_e\,\omega_e+\Delta .
\end{equation}
An \(E\)-surface has all non-zero internal-energy coefficients with the same
orientation after the canonical CFF construction.  An \(H\)-surface is a more
general signed hyperboloid surface.  Pure CFF denominator factors use only
\(E\)-surfaces; LTD and hybrid terms may use \(H\)-surfaces when the residue
geometry requires them.  Numerator-side surfaces are allowed in the numerator
and do not introduce denominator singularities.

\section{Repeated Propagators Without Numerator Derivatives}

Repeated propagators correspond to groups
\begin{equation}
  R_g=\{r_{g,1},\ldots,r_{g,\nu_g}\},\qquad \nu_g\ge2,
\end{equation}
whose members have the same canonical routing and mass, possibly with reversed
edge orientation.  Let \(Q_g^0\) be the canonical repeated-channel energy and
\(\omega_g\) the common on-shell energy.  The purpose of the hybrid
construction is to evaluate the repeated-channel residues without applying
derivatives to the black-box numerator.

\phantomsection
\subsection*{1.2.a The Fully Generic Formula}
\addcontentsline{toc}{subsection}{1.2.a The Fully Generic Formula}

Start from the four-dimensional integral split into non-repeated edges \(S\)
and repeated groups \(G_{\rm rep}\):
\begin{align}
  I_{\Graph,\Basis}[N]
  &=
  \int \prod_{r=1}^{L}\frac{\dd^4 k_r}{(2\pi)^4}\,
  N(\{q_e(k,p)\})
  \prod_{s\in S}D_s^{-1}
  \prod_{g\in G_{\rm rep}}\prod_{a=1}^{\nu_g}D_{r_{g,a}}^{-1}
  \notag\\
  &=
  \int \prod_{r=1}^{L}\frac{\dd^4 k_r}{(2\pi)^4}\,
  N(\{q_e(k,p)\})
  \prod_{s\in S}D_s^{-1}
  \prod_{g\in G_{\rm rep}}D_g^{-\nu_g}.
  \label{eq:collapsed-repeated-integral}
\end{align}
The second equality is only a denominator identity.  The numerator still has
one EMR argument for every original edge \(e\in E\), including all repeated
copies.  Thus the collapsed channel \(g\) has a single denominator energy
\(Q_g^0\), but the numerator still sees \(q_{r_{g,a}}^0\) for
\(a=1,\ldots,\nu_g\).

Let \(\bar E=S\sqcup G_{\rm rep}\) be the collapsed edge set and assign powers
\[
  n_{\bar e}=
  \begin{cases}
    1, & \bar e=s\in S,\\
    \nu_g, & \bar e=g\in G_{\rm rep}.
  \end{cases}
\]
For every rank-\(L\) collapsed cut basis \(B\subset\bar E\) and pole-sign vector
\(\sigma_B\), solve the cut equations
\begin{equation}
  Q_b^0(k^0,p^0)=\sigma_b\omega_b,\qquad b\in B,
  \label{eq:cut-equations}
\end{equation}
for the loop energies:
\begin{equation}
  k^0=K_{B,\sigma_B}(\vec k;p,m).
\end{equation}
The Jacobian associated with this change of variables is
\begin{equation}
  J_{B}=\det(\rho_{br})_{b\in B,\,r=1,\ldots,L}.
\end{equation}
With the contour convention used in the implementation, the sign carried by
the corresponding simple LTD residue is denoted by
\(s_{B,\sigma_B}\in\{\pm1\}\).  The direct residue theorem gives the strict
equality
\begin{align}
  I_{\Graph,\Basis}[N]
  &=
  \int \prod_{r=1}^{L}\frac{\dd^3 k_r}{(2\pi)^3}
  \sum_{B,\sigma_B}
  s_{B,\sigma_B}
  \Biggl[
  \prod_{b\in B}
  \frac{1}{(n_b-1)!}
  \left(\frac{\partial}{\partial Q_b^0}\right)^{n_b-1}
  \notag\\[-2mm]
  &\hspace{34mm}
  \frac{
    N(\{q_e^0(k^0,p^0),\vec q_e\})
  }{
    \prod_{b\in B}(Q_b^0+\sigma_b\omega_b)^{n_b}
    \prod_{\bar e\notin B}D_{\bar e}^{n_{\bar e}}
  }
  \Biggr]_{k^0=K_{B,\sigma_B}} .
  \label{eq:raw-higher-residue}
\end{align}
The determinant \(J_B\) is already included in the LTD residue sign and
normalization used to define \(s_{B,\sigma_B}\).  For the DOT examples all
non-zero \(J_B\) are \(\pm1\).  If a different integer LMB gives
\(|J_B|\ne1\), the residue weight is multiplied by \(1/|J_B|\).

Equation~\eqref{eq:raw-higher-residue} is not the final representation when a powered
channel is cut, because the derivatives act on \(N\).  Define the set of active
original edges
\begin{equation}
  \Active_B
  =
  \left\{
  e\in E\;:\;
  q_e^0 \hbox{ depends on } Q_b^0
  \hbox{ for some } b\in B \hbox{ with } n_b>1
  \right\}.
  \label{eq:active-edges}
\end{equation}
For the affine EMR class, one exact algebraic route is to insert the
multilinear pole interpolation identity
\begin{equation}
  N(\{q_e^0\})
  =
  \sum_{\theta\in\{\pm1\}^{\Active_B}}
  N\!\left(\{q_{e,\theta}^0\}\right)
  \prod_{e\in\Active_B}
  \frac{1+\theta_e q_e^0/\omega_e}{2},
  \label{eq:affine-interpolation}
\end{equation}
with \(q_{e,\theta}^0=\theta_e\omega_e\) for active edges and the ordinary
affine value for inactive edges.  Since \(N\) is affine in each active EMR
energy, differentiating \eqref{eq:affine-interpolation} differentiates only the
known interpolation factors and never the black-box function values.

The form actually used for the hybrid repeated sectors is the equivalent
physical CFF-localized form.  To derive it independently, temporarily split the
repeated denominator into distinct on-shell energies
\(\omega_{g,a}\), \(a=1,\ldots,\nu_g\), and use the linear-pole
decomposition
\begin{equation}
  \frac{1}{(Q_g^0)^2-\omega_{g,a}^2+\ii0}
  =
  \sum_{\tau_{g,a}=\pm1}
  \frac{\tau_{g,a}}{2\omega_{g,a}}\,
  \frac{1}{Q_g^0-\tau_{g,a}\omega_{g,a}+\ii0\,\tau_{g,a}}.
  \label{eq:linear-pole-decomposition}
\end{equation}
Each linear pole has the causal exponential representation
\begin{equation}
  \frac{1}{Q_g^0-\tau\omega+\ii0\,\tau}
  =
  -\ii\,\tau\int_0^\infty \dd \lambda\,
  \exp\!\left[
    \ii\tau\lambda Q_g^0-\ii\lambda\omega
  \right] .
  \label{eq:linear-pole-exponential}
\end{equation}
For a repeated group this gives a sign orbit
\(\tau_g=(\tau_{g,1},\ldots,\tau_{g,\nu_g})\) and the phase
\begin{equation}
  \Phi_{\tau_g}
  =
  \exp\!\left[
    \ii Q_g^0\sum_{a=1}^{\nu_g}\tau_{g,a}\lambda_{g,a}
    -\ii\sum_{a=1}^{\nu_g}\lambda_{g,a}\omega_{g,a}
  \right],
  \qquad \lambda_{g,a}\ge0 .
  \label{eq:repeated-cff-phase}
\end{equation}
The loop-energy integrations are now elementary Fourier integrations.  For a
complete set of CFF/LTD pole choices \(P\), with signs \(\tau_j\), energies
\(\omega_j\), and loop-energy rows \(\rho_j\), one obtains the localizing
constraint
\begin{equation}
  \prod_{r=1}^{L}
  2\pi\,
  \delta\!\left(
    \sum_{j\in P}\tau_j\lambda_j\rho_{jr}
  \right).
  \label{eq:tau-localizing-delta}
\end{equation}
Solving \eqref{eq:tau-localizing-delta} with \(\lambda_j\ge0\) is the
standard CFF source/sink contraction step: the surviving sign assignments are
precisely the causal orientations.  In a repeated group the individual
\(\tau_{g,a}\) survive as numerator-sampling signs, while the localized
collapsed chain has a single sign \(\chi_g=\pm1\) for the energy \(Q_g^0\).
If all repeated-copy signs are equal, \(\chi_g\) is forced to that common sign.
If the orbit is mixed, both \(\chi_g=\pm1\) are possible localizations.  This
is exactly the rule implemented as
\[
  \tau_g=(+,+,+)\Rightarrow \chi_g=+,\qquad
  \tau_g=(-,-,-)\Rightarrow \chi_g=-,\qquad
  \hbox{mixed }\tau_g\Rightarrow \chi_g=\pm .
\]
Only after solving the delta constraints do we take the equal-energy limit
\(\omega_{g,a}\to\omega_g\).  The split expression is an ordinary simple-pole
CFF/LTD expression, hence the limit is equal to the confluent residue
\eqref{eq:raw-higher-residue}.  The difference is organizational:
\begin{itemize}[leftmargin=2em]
  \item the confluent derivative form differentiates the known denominator
        factors and would differentiate \(N\) unless
        \eqref{eq:affine-interpolation} is inserted first;
  \item the physical CFF-localized form samples \(N\) directly at repeated-copy
        on-shell sign orbits \((\tau_g,\chi_g)\), then combines those samples
        with rational affine weights.
\end{itemize}
For an affine EMR numerator the two forms are identical because both are the
same equal-energy limit of the split simple-pole identity.  They differ only in
the basis of numerator calls: the direct interpolation basis uses the formal
active edge signs \(\theta_e\), while the physical basis uses the CFF-localized
orbits \((\tau_g,\chi_g)\).  The latter is the one printed by the current
hybrid JSON labels such as \texttt{physG3tmmmcm}.

It remains to define the rational weights multiplying those physical samples.
Let \(\Samples_{B,\sigma}\) be the finite set of physical samples obtained from
the above sign-orbit rule for the active repeated groups of the cut
\((B,\sigma_B)\).  A sample \(x\in\Samples_{B,\sigma}\) determines complete
loop and edge energy maps,
\[
  K_x^0(\vec k;p,m),\qquad q_{e,x}^0(\vec k;p,m).
\]
Write every component of every map as an affine vector in the symbols
\[
  1,\qquad p_a^0,\qquad \omega_{\bar e},
\]
where repeated copies share the collapsed symbol \(\omega_g\).  Let
\(R_x\) be the column vector containing all those affine coefficients for
sample \(x\).  For a numerator-derivative multi-index \(\beta\), let
\(T_\beta\) be the target affine vector:
\begin{itemize}[leftmargin=2em]
  \item for \(\beta=0\), \(T_0\) is the ordinary cut map
        \(K_{B,\sigma_B}^0,q_{e,B,\sigma_B}^0\);
  \item for \(|\beta|=1\), \(T_\beta\) is the first derivative of those maps
        with respect to the corresponding repeated cut energy, multiplied by
        \(2\omega_g\).  The compensating factor \(1/(2\omega_g)\) is stored as
        one additional repeated half-edge in the JSON.
\end{itemize}
The sample coefficients \(\lambda_{\beta x}\) are defined by the affine system
\begin{equation}
  \sum_{x\in\Samples_{B,\sigma}}\lambda_{\beta x}= \delta_{\beta0},
  \qquad
  \sum_{x\in\Samples_{B,\sigma}}\lambda_{\beta x} R_x=T_\beta.
  \label{eq:physical-affine-weights}
\end{equation}
If the system is underdetermined, the implementation uses the minimum Euclidean
norm solution.  Equivalently, after discarding dependent rows, if \(R\) is the
matrix with columns \(R_x\) and \(t_\beta\) is the target vector including the
constant row, then
\begin{equation}
  \lambda_\beta
  =
  R^{T}(RR^{T})^{-1}t_\beta .
  \label{eq:minnorm-affine-weights}
\end{equation}
This equation is not a fit or approximation: the affine support of
\(\Samples_{B,\sigma}\) is checked exactly over rational numbers.

Substituting \eqref{eq:affine-interpolation} into
\eqref{eq:raw-higher-residue} gives
\begin{align}
  I_{\Graph,\Basis}[N]
  &=
  \int \prod_{r=1}^{L}\frac{\dd^3 k_r}{(2\pi)^3}\,
  I^{\rm HybridLTD}_{\Graph,\Basis,\kappa}[N](\vec k;p,m),
  \label{eq:spatial-hybrid-integral}\\
  I^{\rm HybridLTD}_{\Graph,\Basis,\kappa}[N]
  &=
  \sum_{\mathfrak o\in\Orient_{\rm H}(\Graph,\Basis,\kappa)}
  N\!\left(Q_{\mathfrak o}^{\rm H}(\vec k;p,m)\right)
  W_{\mathfrak o}^{\rm H}(\vec k;p,m).
  \label{eq:hybrid-master}
\end{align}
Here one orientation is
\[
  \mathfrak o=(B,\sigma_B,x,\beta,\gamma,\mu),
\]
where \(x\in\Samples_{B,\sigma}\), \(\beta\) is the part of the confluent
derivative reconstructed from numerator samples, \(\gamma=\alpha-\beta\) acts
on known denominator factors, and \(\mu\) labels one algebraic term in the
known denominator derivative expansion.  Here
\[
  \alpha_b=n_b-1,\qquad b\in B.
\]
The numerator is never differentiated; it is only sampled at
\[
  Q_{\mathfrak o}^{\rm H}=\{q_{e,x}^0,\vec q_e\}_{e\in E}.
\]

The denominator derivative coefficient is also explicit.  Each known factor is
of the form
\[
  F_j(y)=\zeta_j(y)^{-p_j},\qquad
  \zeta_j(y)=\zeta_j(0)+a_j\cdot y,
\]
where \(y=(Q_b^0-\sigma_b\omega_b)_{b\in B}\), \(p_j\) is either a propagator
power or a surface power, and \(a_j\) is the vector of derivatives with respect
to the cut-basis energies.  A multi-derivative assignment
\(\delta_j\) with \(\sum_j\delta_j=\gamma\) contributes
\begin{equation}
  d_{\{\delta_j\}}
  =
  \frac{\gamma!}{\prod_j\delta_j!}
  \prod_j
  (-1)^{|\delta_j|}
  (p_j)_{|\delta_j|}
  a_j^{\delta_j},
  \qquad
  (p)_r=p(p+1)\cdots(p+r-1).
  \label{eq:denominator-derivative-coeff}
\end{equation}
For multi-indices,
\(\gamma!=\prod_b\gamma_b!\),
\(\delta_j!=\prod_b(\delta_j)_b!\), and
\[
  a_j^{\delta_j}=\prod_b (a_j)_b^{(\delta_j)_b}.
\]
For a surface factor \(\zeta_j=\eta_j\), the derivative raises the denominator
power of that cached surface.  For a cut half-edge factor
\((2\omega_b)^{-p_j}\), it raises the number of stored half-edge entries.

The complete rational coefficient multiplying one JSON variant is therefore
\begin{equation}
  c_{\mathfrak o}
  =
  s_{B,\sigma_B}\,
  \frac{1}{\alpha!}\,
  \binom{\alpha}{\beta}\,
  \lambda_{\beta x}\,
  d_{\mu},
  \label{eq:hybrid-coefficient}
\end{equation}
where \(d_\mu\) is the value of
\eqref{eq:denominator-derivative-coeff} for the derivative assignment
\(\mu=\{\delta_j\}\), and
\[
  \alpha!=\prod_{b\in B}\alpha_b!,\qquad
  \binom{\alpha}{\beta}=\prod_{b\in B}\binom{\alpha_b}{\beta_b}.
\]
The weight is the finite product
\begin{equation}
  W_{\mathfrak o}^{\rm H}
  =
  c_{\mathfrak o}
  \prod_{h\in H_{\mathfrak o}}\frac{1}{2\omega_h}
  \prod_{\eta\in D_{\mathfrak o}}\eta(\vec k;p,m)^{-m_{\mathfrak o,\eta}}
  \prod_{\zeta\in N_{\mathfrak o}}\zeta(\vec k;p,m) .
  \label{eq:hybrid-weight}
\end{equation}
Here \(H_{\mathfrak o}\) is the multiset of half-edge ids, so each
\(h\in H_{\mathfrak o}\) contributes \((2\omega_h)^{-1}\).  The denominator
multiset \(D_{\mathfrak o}\) contains cached surfaces
\[
  \eta(\vec k;p,m)=\sum_e c_e\omega_e+\sum_a d_a p_a^0+c_0,
\]
with \(c_e,d_a,c_0\in\mathbb Z\).  The numerator multiset
\(N_{\mathfrak o}\) contains cached numerator-only surfaces of the same
linear form; each \(\zeta\in N_{\mathfrak o}\) is exactly a known linear
factor in on-shell and external energies appearing in the numerator, not an
additional denominator.  If no repeated channel is
present, then
\(n_b=1\) for all cuts, \(\Active_B=\varnothing\), and
\eqref{eq:hybrid-master} collapses exactly to the LTD representation with the
same \((\Graph,\Basis,\kappa)\).

\phantomsection
\subsection*{1.2.b The Box with a Triple Repeated Propagator}
\addcontentsline{toc}{subsection}{1.2.b The Box with a Triple Repeated Propagator}

Consider the one-loop box routing
\[
  q_0=k,\qquad q_1=k+p_1,\qquad q_2=k+p_1+p_2,\qquad
  Q=k+p_1+p_2+p_3,
\]
with three repeated copies of \(Q\).  The collapsed denominator is
\[
  D_0D_1D_2D_Q^3 .
\]
Write \(\Delta_i=p_1^0+\cdots+p_i^0\) and define the causal surfaces
\begin{align}
  E_{ij}^{\sigma_i\sigma_j}
  &=
  \sigma_i\omega_i+\sigma_j\omega_j+\Delta_i-\Delta_j,\\
  E_{iQ}^{\sigma_i\sigma_Q}
  &=
  \sigma_i\omega_i+\sigma_Q\omega_Q+\Delta_i-\Delta_3 .
\end{align}
The three simple-cut sectors are
\begin{equation}
  I_i^{\rm simple}
  =
  -\frac{
    N\!\left(
      q_0^0,q_1^0,q_2^0,Q^0,Q^0,Q^0
    \right)_{k^0=-\Delta_i+\omega_i}
  }{
    2\omega_i\,
    E_{ij}^{++}\,
    E_{i\ell}^{++}\,
    \left(E_{iQ}^{++}\right)^3
  },
  \qquad \{i,j,\ell\}=\{0,1,2\}.
  \label{eq:box-simple-cuts}
\end{equation}
The signs shown correspond to the upper on-shell pole convention used in the
examples; the opposite closure is obtained by replacing every displayed
\(+\omega\) by the corresponding signed on-shell energy.

For the repeated cut used in the printed example, the collapsed LTD cut is
\(Q^0=+\omega_Q\), hence
\[
  k^0=\omega_Q-\Delta_3 .
\]
The six denominator surfaces generated by the three non-cut edges are
\begin{align}
  \eta_{18}&=-\omega_0+\omega_Q-\Delta_3,&
  \eta_{19}&=\phantom{-}\omega_0+\omega_Q-\Delta_3,\\
  \eta_{20}&=-\omega_1+\omega_Q-\Delta_2,&
  \eta_{21}&=\phantom{-}\omega_1+\omega_Q-\Delta_2,\\
  \eta_{22}&=-\omega_2+\omega_Q-\Delta_1,&
  \eta_{23}&=\phantom{-}\omega_2+\omega_Q-\Delta_1 .
  \label{eq:box-b3-surfaces}
\end{align}
These are precisely surface ids \(18,\ldots,23\) in the CLI output.  The
overview table prints a compact subset of denominator ids, but
\texttt{--show-details-for-orientation} displays the full factorization tree
containing all six surfaces.

Let \(\tau=(\tau_3,\tau_4,\tau_5)\) be the physical repeated-copy signs and
\(\chi\) the localized collapsed-chain sign.  A physical numerator call is
\begin{equation}
  N_{\tau,\chi}
  =
  N\!\left(
    \chi\omega_Q-\Delta_3,\;
    \chi\omega_Q-\Delta_2,\;
    \chi\omega_Q-\Delta_1,\;
    \tau_3\omega_Q,\;
    \tau_4\omega_Q,\;
    \tau_5\omega_Q
  \right).
  \label{eq:box-repeated-samples}
\end{equation}
The admissible physical sample set is
\begin{equation}
  \Samples_Q=
  \{(+++, +), (---,-)\}
  \cup
  \{(\tau,+),(\tau,-):\tau\in\{\pm\}^3,\ \tau\ne+++,\ \tau\ne---\}.
  \label{eq:box-physical-sample-set}
\end{equation}
For each \(x=(\tau,\chi)\in\Samples_Q\), the affine coefficients
\(\lambda_{0x}\) and \(\lambda_{1x}\) defined by
\eqref{eq:physical-affine-weights} reconstruct respectively the numerator
itself and the first numerator derivative.  The full table is
\begin{center}
\begin{tabular}{lrr}
\toprule
sample \(x=(\tau,\chi)\) & \(\lambda_{0x}\) & \(\lambda_{1x}\)\\
\midrule
\((+++,+)\) & \(43/112\) & \(5/8\)\\
\((---,-)\) & \(-27/112\) & \(-5/8\)\\
\((--+,-)\) & \(-3/56\) & \(-1/4\)\\
\((--+,+)\) & \(1/112\) & \(-1/8\)\\
\((-+-,-)\) & \(-3/56\) & \(-1/4\)\\
\((-+-,+)\) & \(1/112\) & \(-1/8\)\\
\((-++,-)\) & \(15/112\) & \(1/8\)\\
\((-++,+)\) & \(11/56\) & \(1/4\)\\
\((+--,-)\) & \(-3/56\) & \(-1/4\)\\
\((+--,+)\) & \(1/112\) & \(-1/8\)\\
\((+-+,-)\) & \(15/112\) & \(1/8\)\\
\((+-+,+)\) & \(11/56\) & \(1/4\)\\
\((++-,-)\) & \(15/112\) & \(1/8\)\\
\((++-,+)\) & \(11/56\) & \(1/4\)\\
\bottomrule
\end{tabular}
\end{center}
The first column sums to one and reconstructs the base affine numerator map.
The second column sums to zero and reconstructs the derivative map, with one
extra factor \(1/(2\omega_Q)\).

For the repeated box, \(\alpha=2\).  Hence \(\beta=0,\gamma=2\) or
\(\beta=1,\gamma=1\).  Define
\[
  P=\prod_{i=18}^{23}\eta_i^{-1},\qquad
  h_r=\sum_{18\le i_1\le\cdots\le i_r\le 23}
       \eta_{i_1}^{-1}\cdots\eta_{i_r}^{-1},\qquad h_0=1 .
\]
The denominator derivative terms are
\begin{align}
  \gamma=2:\quad&
  2\,\frac{P h_2}{(2\omega_Q)^3}
  +6\,\frac{P h_1}{(2\omega_Q)^4}
  +12\,\frac{P}{(2\omega_Q)^5},
  \label{eq:box-den-gamma2}\\
  \gamma=1:\quad&
  -\,\frac{P h_1}{(2\omega_Q)^3}
  -3\,\frac{P}{(2\omega_Q)^4}.
  \label{eq:box-den-gamma1}
\end{align}
These equations are exactly the two factorization trees printed in the JSON:
\(P h_2\) is the large tree with all weakly ordered pairs of surfaces,
\(P h_1\) is the smaller tree, and \(P\) is the single chain
\([18,19,20,21,22,23]\).

The complete repeated-cut contribution is therefore
\begin{equation}
  I_Q^{\rm rep}
  =
  \sum_{x\in\Samples_Q}
  N_x
  \left[
    -\frac{\lambda_{0x}}{2}
    \left(
      2\,\frac{P h_2}{(2\omega_Q)^3}
      +6\,\frac{P h_1}{(2\omega_Q)^4}
      +12\,\frac{P}{(2\omega_Q)^5}
    \right)
    -\lambda_{1x}
    \left(
      -\,\frac{P h_1}{(2\omega_Q)^4}
      -3\,\frac{P}{(2\omega_Q)^5}
    \right)
  \right].
  \label{eq:box-repeated-explicit}
\end{equation}
The overall minus sign is the one-loop LTD residue sign for the \(Q\)-cut in
this routing.  Equation~\eqref{eq:box-repeated-explicit} contains no undefined
coefficient: every rational number is either one of the affine weights in the
table or one of the derivative factors in
\eqref{eq:box-den-gamma2}--\eqref{eq:box-den-gamma1}.

For the concrete orientation requested in the CLI output,
\[
  x=(---,-),\qquad
  N_x=N_{---,-}.
\]
The coefficients are
\[
  \lambda_{0x}=-\frac{27}{112},
  \qquad
  \lambda_{1x}=-\frac{5}{8}.
\]
Using \eqref{eq:hybrid-coefficient}, the \(\beta=0,\gamma=2\) variants are
\[
  -\frac{1}{2}\lambda_{0x}\,(2,6,12)
  =
  \left(\frac{27}{112},\frac{81}{112},\frac{81}{56}\right),
\]
with half-edge powers \(3,4,5\).  The \(\beta=1,\gamma=1\) variants include the
extra numerator-derivative half-edge and have
\[
  -\lambda_{1x}\,(-1,-3)
  =
  \left(-\frac{5}{8},-\frac{15}{8}\right),
\]
with half-edge powers \(4,5\).  These five numbers are exactly
\[
\begin{array}{c|c|c}
\hbox{variant label} & \hbox{prefactor} & \hbox{half-edge multiset}\\
\hline
\texttt{B3|b0|g2|n0:physG3tmmmcm} & 27/112 & [3,3,3]\\
\texttt{B3|b0|g2|n0:physG3tmmmcm} & 81/112 & [3,3,3,3]\\
\texttt{B3|b0|g2|n0:physG3tmmmcm} & 81/56 & [3,3,3,3,3]\\
\texttt{B3|b1|g1|nD0phys:physG3tmmmcm} & -5/8 & [3,3,3,3]\\
\texttt{B3|b1|g1|nD0phys:physG3tmmmcm} & -15/8 & [3,3,3,3,3].
\end{array}
\]
The label \texttt{tmmmcm} means
\(\tau=(-,-,-)\) and \(\chi=-\); \texttt{b0/g2} means
\(\beta=0,\gamma=2\), and \texttt{b1/g1} means
\(\beta=1,\gamma=1\).  The repeated edge signs in the orientation are all
minus because the numerator sample uses the physical repeated-copy signs
\(\tau=(-,-,-)\), even though the denominator tree is the derivative expansion
around the collapsed \(Q^0=+\omega_Q\) cut.

The final boxed formula is
\[
  I^{\rm HybridLTD}_{\Box^{(3)}} =
  I_0^{\rm simple}+I_1^{\rm simple}+I_2^{\rm simple}+I_Q^{\rm rep}.
\]

\phantomsection
\subsection*{1.2.c Sunrise with One Repeated Dot Numerator}
\addcontentsline{toc}{subsection}{1.2.c Sunrise with One Repeated Dot Numerator}

The file \texttt{examples/graphs/sunrise\_pow4.dot} gives a compact multiloop
example where the same repeated-channel mechanism is used in a two-loop LMB.
Its internal EMR momenta are
\[
  q_0=-k_1-k_2+p_1,\qquad q_1=k_1,\qquad
  q_2=q_3=q_4=q_5=k_2 ,
\]
with masses \(m_1,m_2,m_3,m_3,m_3,m_3\).  The repeated group is
\[
  R=\{2,3,4,5\},\qquad \nu_R=4,\qquad Q_R^0=k_2^0,\qquad
  \omega_R=\omega_2=\omega_3=\omega_4=\omega_5 .
\]
Consider the numerator
\[
  N(q)=q_2\cdot p_1
  =
  q_2^0p_1^0-\vec q_2\cdot\vec p_1 ,
  \label{eq:sunrise-one-dot-numerator}
\]
which is affine in the repeated energy.  The command-line shorthand
\[
\begin{aligned}
&\texttt{--energy-degree-bounds 2:1 --numerator-expr auto}\\
&\hspace{8mm}\Longrightarrow\quad
  \texttt{dot(edges[2], ext[0])}
\end{aligned}
\]
generates exactly \eqref{eq:sunrise-one-dot-numerator}.

There are two types of rank-two cut bases in the collapsed graph.  The first
type cuts the two non-repeated edges \(\{0,1\}\).  Since the repeated channel
is not cut, \(n_0=n_1=1\), \(\alpha=0\), and the hybrid term is just the LTD
term with the ordinary numerator map.  In the pretty output this is the
orientation with variant
\[
  \texttt{B0,1|b00|g00|n0:phys0}.
\]

The second type cuts one non-repeated edge and the repeated channel,
\[
  B=\{a,R\},\qquad a\in\{0,1\}.
\]
For definiteness take \(B=\{0,R\}\).  The cut equations are
\[
  -k_1^0-k_2^0+p_1^0=\sigma_0\omega_0,\qquad
  k_2^0=\chi\,\omega_R ,
\]
so
\[
  k_2^0=\chi\omega_R,\qquad
  k_1^0=p_1^0-\sigma_0\omega_0-\chi\omega_R .
  \label{eq:sunrise-cut-map}
\]
The repeated physical sample has
\[
  \tau_R=(\tau_2,\tau_3,\tau_4,\tau_5)\in\{\pm1\}^4 ,
\]
and the CFF-localized sign \(\chi\) obeys the same rule as before: all equal
orbits force \(\chi\) to the common sign, while mixed orbits allow both
\(\chi=\pm1\).  The numerator call for
\eqref{eq:sunrise-one-dot-numerator} is therefore
\begin{equation}
  N_{\tau,\chi}^{\rm sun}
  =
  \tau_2\,\omega_R\,p_1^0-\vec k_2\cdot\vec p_1 .
  \label{eq:sunrise-physical-sample}
\end{equation}
The other repeated signs \(\tau_3,\tau_4,\tau_5\) do not enter this particular
dot numerator, but they are still part of the representation-level numerator
energy map stored in the JSON.  This is deliberate: the JSON describes a
representation valid for any numerator respecting the bound, not only for the
single displayed dot.

Because \(\nu_R=4\), the confluent derivative order is
\[
  \alpha_R=\nu_R-1=3.
\]
For an affine repeated numerator, only \(\beta_R=0\) and \(\beta_R=1\) appear.
Consequently the repeated-cut contribution for any fixed non-repeated cut
sign \(\sigma_0\) has the form
\begin{align}
  I^{\rm sun}_{0,R}
  &=
  \sum_{x=(\tau,\chi)\in\Samples_R}
  N_x^{\rm sun}
  \sum_{\beta=0}^{1}
  \sum_{\mu:\,|\mu|=3-\beta}
  s_{0,R,\sigma_0,\chi}\,
  \frac{1}{3!}\binom{3}{\beta}\,
  \lambda_{\beta x}\,
  d_\mu\,
  \prod_{h\in H_{\beta,\mu}}\frac{1}{2\omega_h}
  \prod_{\eta\in D_{\mu}}\eta^{-1}.
  \label{eq:sunrise-repeated-master}
\end{align}
Here \(\Samples_R\) is the four-copy version of
\eqref{eq:box-physical-sample-set}, \(\lambda_{\beta x}\) is obtained from the
same rational affine system \eqref{eq:physical-affine-weights}, and \(d_\mu\)
is the explicit denominator-derivative coefficient of
\eqref{eq:denominator-derivative-coeff}.  For \(\beta=0\) the half-edge
multiset contains four copies of the repeated edge before additional
derivatives are distributed; for \(\beta=1\) it contains five copies because
the numerator derivative reconstruction stores one extra factor
\((2\omega_R)^{-1}\).

In the actual pretty table the first repeated-channel rows have labels such as
\[
  \texttt{-x++++}\qquad
  \texttt{B0,2|b00|g03|n0:physG2tppppcp}.
\]
This means:
\[
\begin{array}{ll}
\texttt{-x++++}:&
q_0^0=-\omega_0,\ q_1^0 \hbox{ is the solved affine map of }
\eqref{eq:sunrise-cut-map},\
q_2^0=q_3^0=q_4^0=q_5^0=+\omega_R,\\[1mm]
\texttt{B0,2}:&\hbox{the collapsed cut basis is edge }0
\hbox{ and the repeated representative }2,\\
\texttt{b00|g03}:&\beta=(0,0),\ \gamma=(0,3),\\
\texttt{physG2tppppcp}:&
\tau_R=(+,+,+,+),\ \chi=+ .
\end{array}
\]
For a mixed sign row such as \texttt{physG2tmmmpcm}, the four letters after
\texttt{t} are \(\tau_2,\tau_3,\tau_4,\tau_5\), and the final
\texttt{c m} gives \(\chi=-\).  Substituting those signs into
\eqref{eq:sunrise-physical-sample}, and the corresponding
\(\lambda_{\beta x}\) into \eqref{eq:sunrise-repeated-master}, reproduces the
prefactors and half-edge powers printed in the JSON.

\section{Higher Energy Powers in Numerators}

The CFF construction is naturally compatible with numerators affine in each
EMR edge energy.  When a user supplies a bound \(d_e\) for the highest power of
each edge energy in \(N\), the implementation first checks loop-energy UV
convergence.  In every loop-energy direction \(v\), the degree of the numerator
minus twice the number of denominators depending on that direction must be
strictly smaller than \(-1\).  If this check fails, finite pole data are not
enough and the command refuses to build a representation.

\phantomsection
\subsection*{1.3.1.a Quadratic Bounds}
\addcontentsline{toc}{subsection}{1.3.1.a Quadratic Bounds}

Let \(x=q_e^0\), \(\omega=\omega_e\), and suppose \(N\) has degree at most two
in \(x\), while all other EMR edge energies are held fixed.  There is an exact
polynomial division
\begin{equation}
  N(x)=R_eN(x)+\left(x^2-\omega^2\right)C_eN,
  \label{eq:quadratic-division}
\end{equation}
where the pole-safe affine remainder is
\begin{equation}
  R_eN(x)=
  \frac{x+\omega}{2\omega}N(+\omega)
  +
  \frac{\omega-x}{2\omega}N(-\omega),
  \label{eq:quadratic-remainder}
\end{equation}
and the contact quotient is
\begin{equation}
  C_eN=
  \frac{N(+\omega)-2N(0)+N(-\omega)}{2\omega^2}.
  \label{eq:quadratic-contact}
\end{equation}
Therefore
\begin{equation}
  \frac{N}{D_e\prod_{f\ne e}D_f}
  =
  \frac{R_eN}{D_e\prod_{f\ne e}D_f}
  +
  \frac{C_eN}{\prod_{f\ne e}D_f}.
  \label{eq:quadratic-master}
\end{equation}
The first term is affine in \(q_e^0\) and can be handled by ordinary CFF.  The
second term is a pinched graph with edge \(e\) removed, also handled by CFF.
Consequently the pure CFF denominator remains \(E\)-surface only.  The
non-trivial part is to lift the CFF expression of the deleted graph back to an
orientation of the original graph.  The lift is completely mechanical:
\begin{enumerate}[leftmargin=2em]
  \item Delete the pinched edge set \(P\) from the denominator graph and build
        a CFF representation of the lower graph
        \(\Graph\setminus P\).  If the lower graph factorizes into independent
        loop-energy components, build one CFF factor per component and multiply
        the component representations.
  \item Every lower-sector CFF term fixes loop energies
        \(K_{\mathfrak t}^{0,\,\Graph\setminus P}\) in the original LMB.  If
        component CFF factors were built separately, the global loop-energy map
        is reconstructed by solving the full-rank linear system defined by the
        union of component basis edges.
  \item For every original edge, including each deleted edge, compute the
        lifted EMR energy
        \begin{equation}
          \widehat q_{e,\mathfrak t}^{\,0}
          =
          \rho_e\cdot K_{\mathfrak t}^{0,\,\Graph\setminus P}
          +\eta_e\cdot p^0 .
          \label{eq:lifted-deleted-edge-energy}
        \end{equation}
        This number is used only for known polynomial factors.  It is not
        automatically the numerator sample value of a deleted edge.
  \item The numerator sample value of a deleted quadratic edge is one of
        \(+\omega_e,0,-\omega_e\), because
        \(C_eN\) in \eqref{eq:quadratic-contact} depends on those three
        black-box samples.  The JSON orientation therefore assigns
        \(q_e^0=+\omega_e\), \(0\), or \(-\omega_e\) to the deleted edge and
        records the corresponding orientation character \(+\), \(0\), or \(-\).
  \item Any lifted factor \(\widehat q_{e,\mathfrak t}^{\,0}\) that remains
        outside the black-box numerator is stored as a numerator-side surface.
        In the strictly quadratic contact quotient no such \(q\)-factor is
        needed; for cubic and higher bounds it is needed.
  \item The lower-sector CFF denominator surfaces are copied into the original
        surface cache after replacing every lower-sector edge id by its
        original edge id.  No denominator surface involving a deleted edge is
        introduced.
\end{enumerate}
The rational coefficients in the JSON follow directly from
\eqref{eq:quadratic-contact}.  Since half-edge lists encode
\((2\omega_e)^{-1}\), the factor \(1/(2\omega_e^2)\) is represented as
\[
  \frac{1}{2\omega_e^2}
  =
  \frac{2}{(2\omega_e)^2}.
\]
The three samples \(N(+\omega_e),N(0),N(-\omega_e)\) therefore carry raw
half-edge prefactors \(2,-4,2\).  A one-edge CFF pinch contributes one
additional sign \((-1)\) in the implementation's CFF orientation convention, so
the printed pinched quadratic prefactors are \(-2,+4,-2\).

In hybrid mode, the same division is applied to non-affine edge energies before
or inside the repeated-channel residue expansion.  For non-repeated edges this
is identical to the pure CFF lift above.  For repeated edges the division is
combined with the repeated-channel sample set
\(\Samples_{B,\sigma}\).  The black-box numerator is still sampled only at the
explicit finite set of energies supplied by the JSON orientation map.

\phantomsection
\subsection*{1.3.1.b One-Loop Box with One Quadratic Edge}
\addcontentsline{toc}{subsection}{1.3.1.b One-Loop Box with One Quadratic Edge}

For a non-repeated one-loop box with denominators \(D_0D_1D_2D_3\), assume the
only non-affine dependence is \(d_0=2\).  Use the routing
\[
  q_0=k,\quad q_1=k+p_1,\quad q_2=k+p_1+p_2,\quad
  q_3=k+p_1+p_2+p_3,
\]
and set \(\Delta_i=p_1^0+\cdots+p_i^0\).  Equations
\eqref{eq:quadratic-division}--\eqref{eq:quadratic-master} give
\begin{equation}
  I_{\Box}^{\CFF}[N]
  =
  I_{\CFF}[\Box;R_0N]
  +
  I_{\CFF}[\Box/0;C_0N],
  \label{eq:box-quadratic}
\end{equation}
where \(\Box/0\) is the triangle obtained by deleting propagator \(0\).  More
explicitly,
\begin{align}
  I_{\CFF}[\Box;R_0N]
  &=
  \sum_{\mathfrak o\in\Orient_{\CFF}(\Box)}
  \left[
  \frac{q_0^0(\mathfrak o)+\omega_0}{2\omega_0}N(q_0^0=+\omega_0)
  +
  \frac{\omega_0-q_0^0(\mathfrak o)}{2\omega_0}N(q_0^0=-\omega_0)
  \right]
  W_{\mathfrak o}^{\CFF},\\
  I_{\CFF}[\Box/0;C_0N]
  &=
  \sum_{\mathfrak t\in\Orient_{\CFF}(\Box/0)}
  \frac{
    N(+\omega_0)-2N(0)+N(-\omega_0)
  }{2\omega_0^2}
  W_{\mathfrak t}^{\CFF}.
\end{align}
Both \(W_{\mathfrak o}^{\CFF}\) and \(W_{\mathfrak t}^{\CFF}\) contain only
\(E\)-surface denominators.

The lifted triangle terms make the orientation structure concrete.  If the
triangle CFF term cuts edge \(j\in\{1,2,3\}\) with sign \(\sigma_j\), then
\[
  k^0=\sigma_j\omega_j-\Delta_j,\qquad
  \widehat q_0^{\,0}=\sigma_j\omega_j-\Delta_j .
\]
However, in the contact numerator call the deleted edge energy is not
\(\widehat q_0^{\,0}\); it is one of the three interpolation samples
\[
  q_0^0=+\omega_0,\qquad q_0^0=0,\qquad q_0^0=-\omega_0 .
\]
For example, the lower triangle orientation with original signs
\((-++ )\) on edges \(1,2,3\) contributes three lifted original-box
orientations
\[
  +-++,\qquad 0-++,\qquad --++ ,
\]
where the first character is the assigned numerator value of the deleted edge.
The printed coefficients are
\[
\begin{array}{c|c|c|c}
\hbox{orientation} & \hbox{variant} & \hbox{prefactor} &
       \hbox{half-edge multiset}\\
\hline
+-++ & \texttt{pinch[0]|e0=+} & -2 & [0,0,1,2,3]\\
0-++ & \texttt{pinch[0]|e0=0} & 4 & [0,0,1,2,3]\\
--++ & \texttt{pinch[0]|e0=-} & -2 & [0,0,1,2,3].
\end{array}
\]
The denominator surfaces of that lower triangle are the ordinary CFF
surfaces among edges \(1,2,3\); in the CLI table they appear, for this
orientation, as ids \([4,5]\).  The two extra half-edge entries \(0,0\) encode
the factor \((2\omega_0)^{-2}\).  No denominator singularity associated with
edge \(0\) remains because \(D_0\) has been cancelled by the contact quotient.
This is why the final JSON is still a normal orientation of the original box:
all four numerator edge energies are specified, but only the lower triangle
denominator surfaces are present.

\phantomsection
\subsection*{1.3.2.a Bounds Above Quadratic}
\addcontentsline{toc}{subsection}{1.3.2.a Bounds Above Quadratic}

For \(d_e>2\), the same idea works but the contact quotient is no longer a
constant in the deleted edge energy.  Let \(a_1,\ldots,a_{d_e-1}\) be
non-pole interpolation nodes with \(a_j\ne\pm1\), and write
\[
  u=\frac{q_e^0}{\omega_e}.
\]
The implementation uses the deterministic node list
\[
  0,\;2,\;-2,\;3,\;-3,\ldots
\]
truncated to \(d_e-1\) nodes.  The Lagrange basis polynomial associated with
node \(a_j\) is
\begin{equation}
  L_j(u)=
  \prod_{\ell\ne j}\frac{u-a_\ell}{a_j-a_\ell}.
  \label{eq:lagrange-basis}
\end{equation}
It is characterized by \(L_j(a_i)=\delta_{ij}\).  The pole remainder is still
the affine interpolation at \(u=\pm1\), and the contact quotient is
reconstructed from black-box samples by
\begin{equation}
  C_eN(q_e^0)
  =
  \sum_{j}
  \frac{L_j(q_e^0/\omega_e)}{(a_j^2-1)\omega_e^2}
  \left[
    N(a_j\omega_e)
    -\frac{a_j+1}{2}N(+\omega_e)
    -\frac{1-a_j}{2}N(-\omega_e)
  \right].
  \label{eq:higher-contact}
\end{equation}
The quotient \(C_eN\) is now a known polynomial in the lifted deleted-edge
energy \(\widehat q_{e,\mathfrak t}^{\,0}\) of
\eqref{eq:lifted-deleted-edge-energy}.  Expanding each \(L_j(u)\) gives
\[
  L_j(u)=\sum_{r=0}^{d_e-2}\ell_{jr}u^r .
\]
A monomial contribution is therefore
\begin{equation}
  \ell_{jr}\,
  \frac{(\widehat q_{e,\mathfrak t}^{\,0})^r}
       {(a_j^2-1)\omega_e^{r+2}}
  \left[
    N(a_j\omega_e)
    -\frac{a_j+1}{2}N(+\omega_e)
    -\frac{1-a_j}{2}N(-\omega_e)
  \right].
  \label{eq:higher-contact-monomial}
\end{equation}
In JSON this becomes:
\begin{itemize}[leftmargin=2em]
  \item \(r+2\) entries of edge \(e\) in \texttt{half\_edges}, because
        \(\omega_e^{-(r+2)}=2^{r+2}(2\omega_e)^{-(r+2)}\);
  \item \(r\) numerator-side copies of the lifted surface
        \(\widehat q_{e,\mathfrak t}^{\,0}\), unless that linear form is zero;
  \item one numerator map with \(q_e^0=a_j\omega_e\), \(+\omega_e\), or
        \(-\omega_e\), depending on the black-box sample in the bracket.
\end{itemize}
For several higher-power edges, \eqref{eq:higher-contact-monomial} suggests a
recursive construction: pick an active edge, split its contribution into the
pole remainder and the contact sector, then repeat the operation on the
remaining bounded edges of the same graph or lower graph.  Once a cubic or
higher contact has been taken, the lower graph carries known
polynomial factors
\[
  \prod_s \bigl(\widehat q_{e_s,\mathfrak t}^{\,0}\bigr)^{r_s}.
\]
If a second non-affine edge is then contacted, those known factors and the
black-box numerator samples no longer separate into independent edge-wise
finite-pole divisions.  The completion is obtained by reducing those known
polynomial factors against the remaining quadratic denominators, as described
in Sec.~1.3.2.c below.  Algebraically this is the lower-sector
contact/infinity completion of the finite-pole recursion.  In JSON it is still
encoded with the same ingredients: ordinary \(E\)-surface denominator trees,
half-edge factors, numerator-only linear factors, and one black-box numerator
map per orientation.  The only grammar extension is that a denominator tree may
now contain a unit node, represented by a node with no surfaces and no
children.  Such a node evaluates to \(1\) and represents a terminal tadpole or
unit lower-sector contact.

The implementation exposes the following pure-CFF higher-power cases:
\begin{enumerate}[leftmargin=2em]
  \item arbitrary quadratic-or-lower caps \(d_e\le2\), on one-loop,
        multiloop, split-mass, and repeated-signature graphs, provided the
        loop-energy UV check is satisfied;
  \item arbitrary UV-convergent higher caps on non-repeated edges, using the
        recursive lower-contact completion described below;
  \item arbitrary UV-convergent higher caps on repeated-signature channels,
        using a channel normal form before any individual copy is deleted.
\end{enumerate}
In all accepted pure-CFF cases every denominator tree is built only from
\(E\)-surfaces.  Numerator-side cached factors may be \(E\)- or \(H\)-type
linear forms, but the JSON marks them as numerator-only and the evaluator
multiplies them in the numerator.

The repeated-channel case needs one extra algebraic step.  Let
\[
  C=\{e_1,\ldots,e_\nu\}
\]
be a repeated channel with common canonical energy \(y_C\), on-shell energy
\(\omega_C\), and denominator \(D_C=y_C^2-\omega_C^2\).  If the user-supplied
edge bounds in this channel are \(d_{e_a}\), the black-box numerator is known
to have total channel degree at most
\[
  d_C=\sum_{a=1}^{\nu}d_{e_a}.
\]
Choose the deterministic interpolation nodes
\[
  \Samples_{d_C}=(1,-1,0,2,-2,3,-3,\ldots)
\]
truncated to \(d_C+1\) entries.  The Lagrange polynomial for node \(a\) is
\[
  L_a(u)=\prod_{\substack{b\in\Samples_{d_C}\\ b\ne a}}
  \frac{u-b}{a-b},\qquad u=\frac{y_C}{\omega_C}.
\]
For any black-box numerator, with all channel-copy EMR energies sampled at
their signed values \(q_{e_a}^0=s_{e_a}a\,\omega_{e_a}\), where
\(s_{e_a}=\pm1\) is the sign of the copy relative to the canonical channel,
\[
  N(y_C,\ldots)
  =
  \sum_{a\in\Samples_{d_C}}
  L_a(y_C/\omega_C)\,
  N\!\left(q_{e_a}^0=s_{e_a}a\,\omega_{e_a},\ldots\right).
  \label{eq:channel-interpolation}
\]
Now write
\[
  L_a(u)=\sum_{r=0}^{d_C}\ell_{a,r}u^r,\qquad
  r=2b+p,\quad p\in\{0,1\}.
\]
The normal form of each monomial is
\begin{equation}
  \frac{u^{2b+p}}{D_C^\nu}
  =
  \sum_{j=0}^{b}
  \binom{b}{j}
  \frac{y_C^{p}\,D_C^j}
       {\omega_C^{2j+p}D_C^\nu}.
  \label{eq:channel-normal-form}
\end{equation}
If \(j<\nu\), the term keeps \(\nu-j\) explicit copies of the repeated
denominator and deletes the other \(j\) copies.  In JSON it contributes:
\begin{itemize}[leftmargin=2em]
  \item coefficient \(2^{2j+p}\ell_{a,2b+p}\binom{b}{j}\), because
        \(\omega_C^{-(2j+p)}=2^{2j+p}(2\omega_C)^{-(2j+p)}\);
  \item \(2j+p\) entries of the representative channel edge in
        \texttt{half\_edges};
  \item one numerator-only surface for \(y_C\) if \(p=1\);
  \item a lower CFF denominator graph with exactly \(\nu-j\) copies of the
        channel still present.
\end{itemize}
If \(j\ge\nu\), no denominator copy of \(C\) remains.  The extra
\(D_C^{j-\nu}\) is serialized as numerator-only factors
\[
  (y_C-\omega_C)^{j-\nu}(y_C+\omega_C)^{j-\nu},
\]
again expressed through cached surface keys.  The lower denominator is then the
CFF expression of the graph with the whole repeated channel deleted.  This is
still an \(E\)-surface-only denominator representation; the additional factors
are explicitly numerator-side objects.

This channel normal form is the reason pure CFF can now handle high powers on
repeated propagators without ever invoking LTD contact denominators.  It also
makes the ordinary finite-pole/contact recursion a limiting case: for
\(\nu=1\), \eqref{eq:channel-normal-form} is precisely the division of a
single higher-power edge into a pole remainder and a lower contact sector.

\subsection{Uniform-Scale Repeated-Channel Sampling}

The default interpolation variable in
\eqref{eq:channel-interpolation} is \(u_C=y_C/\omega_C\).  This is compact and
is kept as the default implementation because it produces the smallest
coefficients for the ordinary repeated-channel formula.  However, after
\eqref{eq:channel-normal-form}, high-degree monomials can generate explicit
inverse powers of \(\omega_C\).  These inverse powers are artificial
interpolation factors: the complete orientation sum is independent of the
normalization used for the interpolation nodes.  They are not physical
propagator singularities, and for a massless representative edge they can make
individual terms unnecessarily sensitive to a soft limit before the term sum is
formed.

The implementation therefore also supports a user-provided non-zero scale
\(M\).  Let
\[
  \Samples_d=(1,-1,0,2,-2,\ldots)
\]
be the deterministic integer node list.  In uniform-scale mode the pure CFF
channel interpolant is
\begin{equation}
  N(y_C,\ldots)
  =
  \sum_{a\in\Samples_{d_C}}
  L_a(y_C/M)\,
  N\!\left(q_{e_b}^0=s_{e_b}a\,M,\ldots\right),
  \label{eq:uniform-channel-interpolation}
\end{equation}
where \(s_{e_b}=\pm1\) is the sign of member \(e_b\) relative to the channel
representative.  The Lagrange polynomial is the same polynomial in its
dimensionless argument:
\[
  L_a(v)=\sum_{r=0}^{d_C}\ell_{a,r}v^r,\qquad v=y_C/M.
\]
For a monomial \(r=2b+p\), \(p\in\{0,1\}\), the normal form becomes
\begin{equation}
  \frac{(y_C/M)^{2b+p}}{D_C^\nu}
  =
  \sum_{j=0}^{b}
  \binom{b}{j}
  \frac{y_C^p\,\omega_C^{2(b-j)}\,D_C^j}
       {M^{2b+p}D_C^\nu}.
  \label{eq:uniform-channel-normal-form}
\end{equation}
Compared with \eqref{eq:channel-normal-form}, the inverse scale is now
\(M^{-(2b+p)}\), while positive powers of \(\omega_C\) are numerator-side
known factors.  If \(j<\nu\), the denominator keeps \(\nu-j\) channel copies.
If \(j\ge\nu\), the excess \(D_C^{j-\nu}\) is again serialized as numerator
factors \((y_C-\omega_C)^{j-\nu}(y_C+\omega_C)^{j-\nu}\).  Hence the
denominator surfaces are unchanged CFF causal surfaces; only the numerator map
and variant prefactor gain explicit dependence on \(M\).

For hybrid LTD/CFF, the same idea is applied around the localized LTD pole.
If a repeated basis edge has sign \(\sigma=\pm1\), the finite-difference
samples are
\begin{equation}
  q_C^0=\sigma\,\omega_C+aM,\qquad a\in\Samples_d ,
  \label{eq:uniform-hybrid-samples}
\end{equation}
and the finite-difference weights carry \(M^{-|\beta|}\) for a
multi-derivative index \(\beta\).  This replaces the default
\((2\omega_C)^{-|\beta|}\) factors without changing the final summed value.

The CLI exposes three policies:
\[
\begin{array}{ll}
\texttt{none} & \text{use } \omega_C \text{ normalization everywhere
  (default)},\\
\texttt{beyond-quadratic} & \text{use } M \text{ only when the repeated-channel
  bound is } d_C>2,\\
\texttt{all} & \text{use } M \text{ for every repeated-channel bound }
  d_C>1 .
\end{array}
\]
The second policy keeps the quadratic case in the original compact
\(\omega_C\)-normalized form; the third is useful when one wants all repeated
finite-difference samples to have the same external normalization.  The
runtime evaluator rejects \(M=0\).  Real evaluators accept positive or negative
real \(M\).  Complex-valued Symbolica evaluators also accept complex \(M\);
real-valued Symbolica evaluators reject it.

For the two-loop \texttt{sunrise\_pow4.dot} graph, the repeated channel is
\[
  C=\{2,3,4,5\},\qquad y_C=k_2^0,\qquad \nu=4 .
\]
With
\[
  \texttt{--energy-degree-bounds 2:5}
\]
the channel degree is \(d_C=5\), so
\[
  \Samples_5=(1,-1,0,2,-2,3).
\]
The JSON samples the black-box numerator at the six maps
\[
  q_2^0=a\,\omega_2,\qquad
  q_3^0=a\,\omega_3,\qquad
  q_4^0=a\,\omega_4,\qquad
  q_5^0=a\,\omega_5,
  \qquad a\in\Samples_5,
\]
where \(\omega_2=\omega_3=\omega_4=\omega_5\) numerically because the repeated
copy masses and spatial routings are equal.  For a monomial check,
\[
  \frac{(y_C)^5}{D_C^4}
  =
  \frac{\omega_C^4 y_C}{D_C^4}
  +2\,\frac{\omega_C^2 y_C}{D_C^3}
  +\frac{y_C}{D_C^2}.
  \label{eq:sunrise-channel-monomial-check}
\]
The serialized black-box formula does not assume this monomial; it emits the
six Lagrange sample maps above and the rational coefficients obtained from
\eqref{eq:channel-interpolation}--\eqref{eq:channel-normal-form}.  The three
denominator powers in \eqref{eq:sunrise-channel-monomial-check} correspond to
JSON variants keeping four, three, and two repeated denominator copies,
respectively.  Every such variant has only CFF \(E\)-surfaces in the
denominator and, when \(p=1\), one numerator-only surface representing
\(y_C\).  The tests compare this pure-CFF value with the hybrid value and with
the split-mass LTD limiting proxy.

\paragraph{How the JSON rows are obtained.}
The oriented JSON does not store the polynomial
\eqref{eq:channel-interpolation} symbolically.  It stores the fully expanded
terms of
\[
  L_a(u)D_C^{-\nu}
\]
after applying \eqref{eq:channel-normal-form}.  For each interpolation node
\(a\) and each monomial \(u^{2b+p}\) in \(L_a\), the builder emits one variant
for each \(j=0,\ldots,b\).  The row data are
\[
\begin{array}{rl}
\texttt{sample} &= a,\\
\texttt{remaining\_power} &= \max(0,\nu-j),\\
\texttt{parity} &= p,\\
\texttt{cancelled\_power} &= \max(0,j-\nu),\\
\texttt{inverse\_ose\_power} &= 2j+p,\\
\texttt{coefficient} &= \ell_{a,2b+p}\binom{b}{j}.
\end{array}
\]
The JSON prefactor of that variant is this coefficient multiplied by
\(2^{2j+p}\) and by the ordinary lower CFF prefactor of the denominator graph
with \(\max(0,\nu-j)\) channel copies retained.  The ordinary CFF half-edge
list of that lower graph is then appended with \(2j+p\) copies of the
representative channel edge.  If \(p=1\), the numerator-surface list gets one
cached copy of the canonical channel energy \(y_C\).  This construction is
purely mechanical and is what the evaluator uses directly.

For \texttt{box\_pow3.dot} with
\[
  \texttt{--energy-degree-bounds 3:4},
\]
the repeated channel is \(C=\{3,4,5\}\), \(\nu=3\), \(d_C=4\), and
\(\Samples_4=(1,-1,0,2,-2)\).  The first numerator map printed by
\texttt{build --family cff --pretty} has orientation label
\(\texttt{-+++++}\).  For its sample \(a=1\), the actual variants are:
\[
\begin{array}{c|c|c|c|c|c}
\texttt{pref} & \texttt{remaining\_power} & p &
\texttt{inverse\_ose\_power} & \texttt{half\_edges} &
\texttt{num\_surfaces}\\
\hline
-8/3 & 1 & 0 & 4 & [0,1,2,3,3,3,3,3] & []\\
-4/3 & 2 & 0 & 2 & [0,1,2,3,3,3,4] & []\\
 4/3 & 2 & 1 & 3 & [0,1,2,3,3,3,3,4] & [19]\\
 1/2 & 3 & 0 & 0 & [0,1,2,3,4,5] & []\\
 1   & 3 & 1 & 1 & [0,1,2,3,3,4,5] & [19].
\end{array}
\]
Here surface \(19\) is numerator-only and is the cached linear form for the
canonical repeated energy \(y_C\) in that orientation.  The denominator
surfaces referenced by the factorization trees are all \(E\)-surfaces.

For the mixed box command
\[
  \texttt{--energy-degree-bounds 0:1,1:1,2:0,3:4},
\]
the same channel reduction is performed first.  The affine powers on edges
\(0\) and \(1\) do not require new black-box calls.  They become known linear
factors whenever a lower contact sector leaves them outside the sampled
black-box numerator, so the same orientation label can have additional
variants with extra numerator surfaces, for example surfaces \(20,21,22\) in
the current output.  These variants are still grouped under the same
\texttt{edge\_q0} map; this is the invariant that one orientation equals one
black-box numerator call.

For \texttt{sunrise\_pow4.dot} with
\[
  \texttt{--energy-degree-bounds 2:5},
\]
the repeated channel is \(C=\{2,3,4,5\}\), \(\nu=4\), \(d_C=5\), and
\(\Samples_5=(1,-1,0,2,-2,3)\).  The first numerator map has orientation label
\(\texttt{++++++}\).  For sample \(a=1\), the variants are:
\[
\begin{array}{c|c|c|c|c|c}
\texttt{pref} & \texttt{remaining\_power} & p &
\texttt{inverse\_ose\_power} & \texttt{half\_edges} &
\texttt{num\_surfaces}\\
\hline
-8/3 & 2 & 0 & 4 & [0,1,2,2,2,2,2,3] & []\\
 8/3 & 2 & 1 & 5 & [0,1,2,2,2,2,2,2,3] & [5]\\
-4/3 & 3 & 0 & 2 & [0,1,2,2,2,3,4] & []\\
10/3 & 3 & 1 & 3 & [0,1,2,2,2,2,3,4] & [5]\\
 1/2 & 4 & 0 & 0 & [0,1,2,3,4,5] & []\\
 1   & 4 & 1 & 1 & [0,1,2,2,3,4,5] & [5].
\end{array}
\]
Surface \(5\) is numerator-only and represents the lifted canonical channel
energy
\[
  y_C=-\omega_0-\omega_1+p_1^0
\]
for this orientation.  Its sign-flipped companion appears as surface \(6\) in
the \(\texttt{--++++}\) orientation.  These two surfaces are marked
\((e)\) in pretty output because they occur only in the numerator.

If the same sunrise graph is built with
\[
  \texttt{--energy-degree-bounds 2:3,3:3},
\]
the channel degree is \(d_C=6\), because both powered repeated copies belong
to the same denominator channel.  The node set is
\(\Samples_6=(1,-1,0,2,-2,3,-3)\).  The first printed numerator map has seven
variants for sample \(a=1\); the leading channel-reduction metadata are
\[
\begin{array}{c|c|c|c|c}
\texttt{pref} & \texttt{remaining\_power} & p &
\texttt{inverse\_ose\_power} & \texttt{coefficient}\\
\hline
-4/3  & 1 & 0 & 6 & 1/48\\
-10/3 & 2 & 0 & 4 & -5/24\\
 2/3  & 2 & 1 & 5 & 1/48\\
-13/12& 3 & 0 & 2 & 13/48 .
\end{array}
\]
The remaining three variants are the same expansion continued to
\(\texttt{remaining\_power}=3,4,4\).  The important point is that the JSON
does not distinguish which copy originally carried which cubic cap once the
copies are identified as a repeated channel; it only needs the total channel
degree \(d_C=6\) and the individual numerator sample map for every original
edge copy.

Pure CFF does not need a separate repeated-pole derivative construction:
repeated propagators are repeated denominator factors in the CFF contour
integral.  The channel normal form above simply performs the finite polynomial
division at the level of the repeated denominator power before CFF denominator
trees are generated.  The hybrid formula remains useful for the LTD-like
repeated-channel organization, but it is not needed to make the pure-CFF
higher-power lift well-defined.

Hybrid mode has a different status.  The hybrid repeated-channel formula is
not constrained to keep pure-CFF denominator trees \(E\)-only: \(H\)-surfaces
are allowed because the non-repeated part is LTD-like and CFF is used only to
localize repeated channels.  With \texttt{--energy-degree-bounds}, numerator
derivatives in the repeated-channel residue are replaced by tensor-product
finite-difference samples in the repeated basis coordinates.  For a derivative
multiindex \(\beta\), the degree used along basis direction \(b\) is
\[
  D_b=\sum_e d_e\,\mathbf 1\{\partial q_e^0/\partial z_b\ne0\},
\]
where \(z_b\) is the repeated basis energy.  The finite-difference stencil is
the node set \(0,1,-1,2,-2,\ldots\) truncated to \(D_b+1\) nodes, with weights
fixed by exact polynomial interpolation.  Therefore the numerator is still a
black box: the JSON only asks for samples at explicit EMR energy maps.

The implemented hybrid path is validated for arbitrary user-provided caps
subject to the same loop-energy UV check and practical expression size.  The
test suite covers cubic, quartic, and quintic caps on repeated box, sunrise,
kite, mercedes, and iterated-bubble examples by comparing the hybrid value to
the split-mass LTD limiting proxy.  In the absence of repeated propagators,
hybrid collapses exactly to LTD for any accepted bound vector.

\phantomsection
\subsection*{1.3.2.b One-Loop Box with One Cubic Edge}
\addcontentsline{toc}{subsection}{1.3.2.b One-Loop Box with One Cubic Edge}

For the same one-loop box, take \(d_0=3\).  The implementation chooses the
two non-pole nodes
\[
  a_0=0,\qquad a_1=2.
\]
The Lagrange basis is
\begin{equation}
  L_0(u)=1-\frac{u}{2},\qquad
  L_1(u)=\frac{u}{2}.
  \label{eq:box-cubic-lagrange}
\end{equation}
Substituting these two polynomials into \eqref{eq:higher-contact} gives
\begin{align}
  C_0N(q_0^0)
  &=
  -\frac{1-u/2}{\omega_0^2}
  \left[
    N(0)-\frac12 N(+\omega_0)-\frac12 N(-\omega_0)
  \right]
  \notag\\
  &\quad
  +\frac{u/2}{3\omega_0^2}
  \left[
    N(2\omega_0)-\frac32 N(+\omega_0)+\frac12 N(-\omega_0)
  \right],
  \qquad u=\frac{\widehat q_0^{\,0}}{\omega_0}.
  \label{eq:box-cubic-contact}
\end{align}
For a lower triangle orientation \(\mathfrak t\), the lifted deleted energy is
\[
  \widehat q_{0,\mathfrak t}^{\,0}
  =
  \rho_0\cdot K_{\mathfrak t}^{0,\Box/0}.
\]
For instance, if the triangle cut is edge \(1\) with sign \(-\), then
\[
  \widehat q_{0,\mathfrak t}^{\,0}=-\omega_1-p_1^0,
\]
which appears in the surface cache as a numerator-only \(E\)-surface.  If the
cut is edge \(1\) with sign \(+\), the lifted surface is
\(\omega_1-p_1^0\).  These are the numerator-only surfaces printed as
\((e)\) in the cubic box output.

The monomial coefficients of \eqref{eq:box-cubic-contact} map directly to JSON
variants.  Before the one-edge CFF pinch sign, the finite-pole contact
components for edge \(0\) are
\[
\begin{array}{c|c|c|c}
\hbox{sample }q_0^0 & \hbox{prefactor} &
  \hbox{half-edge multiset} & \hbox{numerator surface}\\
\hline
0 & -4 & [0,0] & \varnothing\\
0 & 4 & [0,0,0] & [\widehat q_0^0]\\
+\omega_0 & 2 & [0,0] & \varnothing\\
+\omega_0 & -4 & [0,0,0] & [\widehat q_0^0]\\
-\omega_0 & 2 & [0,0] & \varnothing\\
-\omega_0 & -4/3 & [0,0,0] & [\widehat q_0^0]\\
2\omega_0 & 4/3 & [0,0,0] & [\widehat q_0^0].
\end{array}
\]
The printed one-loop box CFF table includes the additional pinch sign, so the
visible variants are
\[
\begin{array}{c|c|c|c}
\hbox{orientation label} & \hbox{variant} & \hbox{prefactor}
  & \hbox{meaning}\\
\hline
0-++ & \texttt{pinch[0]|e0=0} & 4
  & -(-4)\,N(q_0^0=0)/(2\omega_0)^2\\
0-++ & \texttt{pinch[0]|e0=0} & -4
  & -(4)\,\widehat q_0^0\,N(q_0^0=0)/(2\omega_0)^3\\
x-++ & \texttt{pinch[0]|e0=2} & -4/3
  & -(4/3)\,\widehat q_0^0\,N(q_0^0=2\omega_0)/(2\omega_0)^3\\
+-++ & \texttt{pinch[0]|e0=+} & -2
  & -(2)\,N(q_0^0=+\omega_0)/(2\omega_0)^2\\
+-++ & \texttt{pinch[0]|e0=+} & 4
  & -(-4)\,\widehat q_0^0\,N(q_0^0=+\omega_0)/(2\omega_0)^3\\
--++ & \texttt{pinch[0]|e0=-} & -2
  & -(2)\,N(q_0^0=-\omega_0)/(2\omega_0)^2\\
--++ & \texttt{pinch[0]|e0=-} & 4/3
  & -(-4/3)\,\widehat q_0^0\,N(q_0^0=-\omega_0)/(2\omega_0)^3 .
\end{array}
\]
The label \(x\) in \(x-++\) indicates that the deleted edge is sampled at the
non-standard finite node \(q_0^0=2\omega_0\), so its detailed energy map must be
read from the orientation details rather than from a single \(+\), \(-\), or
\(0\) symbol.  All denominator surfaces in these rows are the lifted triangle
CFF surfaces; \(\widehat q_0^0\) is a numerator-only surface and cannot create
a denominator singularity.

This cubic box example is the first member of the higher-contact recursion.  If
the other box edges are affine, for example
\((d_0,d_1,d_2,d_3)=(3,1,1,1)\), the formula above is exact and agrees with
LTD.  If one spectator is quadratic, for example \((3,1,2,0)\), a second
contact produces the terminal tadpole derived below.  The implemented command
now emits that completion as an explicit unit denominator node, so the resulting
pure-CFF JSON remains \(E\)-surface-only in the denominator and agrees with LTD
for all UV-convergent one-loop box monomial caps through the convergence range
\(\sum_e d_e\le6\).  The same higher cap is allowed on a graph that contains
repeated propagators, including when the capped edge is one of the repeated
copies, because pure CFF treats those copies as ordinary repeated denominator
factors.

\phantomsection
\subsection*{1.3.2.c Terminal Contact Completion}
\addcontentsline{toc}{subsection}{1.3.2.c Terminal Contact Completion}

The recursive lower-contact completion can terminate on a sector whose
remaining denominator set is smaller than the degree of the known polynomial
factor carried by previous contacts.  The simplest terminal sector is useful
because it can be analyzed completely and because it fixes the unit-tree terms
that appear in the JSON grammar.

Consider the one-loop box with edge energies
\[
  q_e^0=k^0+\Delta_e,\qquad
  D_e=(q_e^0)^2-\omega_e^2,\qquad e=0,1,2,3 ,
\]
where \(\Delta_e\) is the external-energy shift of edge \(e\) and
\(\omega_e=E_e(\vec k)\).  Define
\[
  \delta_{ab}=\Delta_a-\Delta_b ,
  \qquad q_a^0=q_b^0+\delta_{ab}.
\]
One representative terminal sector is
\[
  (d_0,d_1,d_2,d_3)=(0,0,3,3),
  \qquad
  N(q)= (q_2^0)^3(q_3^0)^3 .
\]
For a cubic edge the finite-pole/contact division is just
\begin{equation}
  (q_e^0)^3
  =
  \omega_e^2 q_e^0
  +
  D_e\,q_e^0 .
  \label{eq:cubic-one-edge-division}
\end{equation}
The first term is the pole remainder and is handled by the finite-pole CFF
samples.  The second term cancels \(D_e\) and produces a contact lower sector.
Taking the contact part on both cubic edges produces, among other already
handled terms, the lower bubble
\begin{equation}
  \frac{q_2^0 q_3^0}{D_0D_1}.
  \label{eq:terminal-contact-bubble}
\end{equation}
This bubble is the first terminal sector in which the lower-contact completion
produces a unit denominator tree.
Writing the numerator in terms of the remaining edge \(0\),
\[
  q_2^0=q_0^0+\delta_{20},
  \qquad
  q_3^0=q_0^0+\delta_{30},
\]
gives the exact identity
\begin{align}
  q_2^0q_3^0
  &=
  (q_0^0)^2+(\delta_{20}+\delta_{30})q_0^0
    +\delta_{20}\delta_{30}
  \notag\\
  &=
  D_0+
  \omega_0^2+(\delta_{20}+\delta_{30})q_0^0
    +\delta_{20}\delta_{30}.
  \label{eq:terminal-contact-polynomial-division}
\end{align}
Therefore
\begin{equation}
  \frac{q_2^0q_3^0}{D_0D_1}
  =
  \frac{
    \omega_0^2+(\delta_{20}+\delta_{30})q_0^0
      +\delta_{20}\delta_{30}}
       {D_0D_1}
  +
  \frac{1}{D_1}.
  \label{eq:terminal-contact-tadpole}
\end{equation}
The first term has an affine numerator in the lower bubble and is a standard
CFF object.  The second term is a terminal lower-sector tadpole.  It has no
ordinary \(E\)-surface denominator left after the half-edge factor of edge
\(1\); in the JSON grammar it would be represented as a unit denominator tree
with the half-edge list \([1]\).

With the close-below convention used by the code, this terminal contribution
has the sign of the corresponding CFF branch and evaluates to
\[
  -\,\frac{1}{2\omega_1}
\]
for the sector in \eqref{eq:terminal-contact-bubble}.  The same value is
obtained by comparing the serialized CFF unit-tree term against the LTD
reference at high precision.  For the same random kinematics used in the tests,
\[
  \omega_1=0.788665091489179\ldots,\qquad
  \frac{1}{2\omega_1}=0.633982669444499\ldots .
\]
The related box sectors
\[
  (0,3,0,3),\quad (0,3,3,0),\quad (2,1,0,3),
  \quad (3,1,2,0),
\]
and their permutations show the same pattern: the terminal contribution is a single
terminal tadpole \(\pm1/(2\omega_i)\), with \(i\) determined by the last
remaining denominator in the lower sector.  In the intermediate recursion these
branches are terms with no \(E\)-surface denominator chain.  The schema now
serializes them as a unit denominator node
\[
  \texttt{\{surfaces: [], children: []\}},
\]
so the evaluator multiplies by the half-edge factors and by one.  With this
unit-node encoding, all UV-convergent one-loop box monomial caps satisfying
\(\sum_e d_e\le6\) agree with LTD at arbitrary precision.

The natural one-loop generalization is consequently not another finite-pole
sample, but a lower-sector normal-form recursion.  Let \(S\) be the set of
denominators left after some contacts and let \(P(q^0)\) be the known
polynomial factor multiplying the black-box numerator sample.  Choose
\(a\in S\) and divide \(P\) by
\[
  D_a=(q_a^0)^2-\omega_a^2
\]
after rewriting all \(q_e^0\) in terms of \(q_a^0\):
\begin{equation}
  P(q^0)=D_a\,Q_a(q^0)+R_a(q^0),
  \qquad
  \deg_{q_a^0} R_a\le1 .
  \label{eq:lower-sector-normal-form}
\end{equation}
Then
\begin{equation}
  \frac{P}{\prod_{e\in S}D_e}
  =
  \frac{R_a}{D_a\prod_{e\in S\setminus\{a\}}D_e}
  +
  \frac{Q_a}{\prod_{e\in S\setminus\{a\}}D_e}.
  \label{eq:lower-sector-recursion}
\end{equation}
The first term is again a finite-pole CFF term because its numerator is affine
in the active energy.  The second term is the contact completion and is
recursed on the smaller denominator set.  When one denominator remains,
\eqref{eq:lower-sector-recursion} terminates in a tadpole; when no denominator
remains, UV convergence of the full sector must make the sum of such unit
polynomial branches cancel across contact histories.

The same mechanism extends to multiloop contact sectors whenever the known
polynomial factors reduce to the span of the denominator components that remain.
The smallest diagnostic graph is the two-loop penta-denominator
\[
  D_0(k_1+k_2+p_3)D_1(k_1)D_2(k_1+p_1)D_3(k_2)D_4(k_2+p_2).
\]
For the UV-convergent bounds
\[
  (d_0,d_1,d_2,d_3,d_4)=(1,0,3,0,3),
\]
one contact history sets \(q_2^0=q_4^0=q_0^0=q_1^0=0\) and leaves only
\[
  \frac{(p_1^0)(p_2^0-p_3^0)}{D_3(k_2)} .
\]
The original graph has two loop energies, but this lower sector has rank one.
It is nevertheless not ambiguous: all dependence on the unlocalized \(k_1^0\)
has disappeared into external-energy factors.  The completion is therefore a
single terminal tadpole,
\[
  \frac{p_1^0(p_2^0-p_3^0)}{2\omega_3},
\]
times the half-edge factors accumulated along the contact path.  The code
handles this by decomposing the remaining denominator rows into vector-matroid
components, building a CFF or terminal LTD factor for each non-zero component,
and solving only the localized rank directions.  Free loop-energy directions
are assigned the particular value zero in the derived \(\texttt{loop\_q0}\)
map; the EMR \(\texttt{edge\_q0}\) map is the numerator map and is unchanged by
that bookkeeping choice.

This gives a practical generalization of
\eqref{eq:lower-sector-normal-form}--\eqref{eq:lower-sector-recursion}: recurse
on active higher-power or known-polynomial edges, reduce quotient terms against
the remaining denominators, decompose the lower denominator into matroid
components, and serialize terminal components as unit denominator nodes.  The
pure-CFF tests now include the one-loop box sectors above, repeated-signature
box sectors with quartic caps on repeated copies, high-power repeated sunrise
and kite sectors, two repeated channels in the iterated-bubble topology, and
the two-loop rank-deficient penta sectors; in every accepted case the
denominator surface set is \(E\)-only and the value agrees with LTD or, for
repeated propagators, with the hybrid value and the split-mass LTD limiting
proxy.

\chapter{Python Implementation and JSON Format}

\section{Module Structure}

The implementation is organized around an evaluator-complete intermediate
format.
\begin{itemize}[leftmargin=2em]
  \item \texttt{graph\_io.py} parses DOT graphs, extracts momentum signatures,
        detects repeated channels, and builds split-mass reference graphs.
  \item \texttt{graph\_signatures.py} implements the
        \texttt{graph\_from\_signatures} command, which turns products such as
        \texttt{prop(k1+p1,m)} into DOT graphs.
  \item \texttt{orientation\_bundle.py} constructs LTD, CFF, hybrid, and
        bounded-degree contact bundles.  This is where repeated channels are
        collapsed, active numerator maps are enumerated, and CFF lower minors
        are generated.
  \item \texttt{structure.py} serializes bundles to schema-v6 JSON and contains
        the arbitrary-precision Python evaluator.
  \item \texttt{pretty.py} renders the JSON in a human-readable form.
  \item \texttt{symbolica\_eval.py} compiles a fixed-numerator JSON expression
        to Symbolica eager and assembly-backed evaluators.
\end{itemize}

\section{Bounded-Degree Support Matrix}

The command-line option \texttt{--energy-degree-bounds} is interpreted as an
upper bound \(d_e\) for the degree of the black-box numerator in each EMR edge
energy \(q_e^0\).  The builder first applies the loop-energy UV check described
in Section~1.3.  If the check fails, no representation is built.
For the \texttt{test} and \texttt{compare} subcommands this option is common:
the same bound vector is used to build the CFF expression, the hybrid
expression, and the split-mass LTD reference.  The
\texttt{--cff-energy-degree-bounds} option is retained only for CFF-only
diagnostics.

When a command both evaluates a numerator and sees
\texttt{--numerator-expr auto}, the CLI replaces \texttt{auto} by a concrete
black-box numerator before any representation is evaluated or compiled.  The
\texttt{build} command accepts the same option as metadata: it records the
resolved expression in \texttt{graph.numerator\_expr} while leaving the
orientation structure unchanged.  If the graph has \(n_{\rm ext}\) external
momentum symbols and the normalized bound vector is
\(d=(d_0,\ldots,d_{|E|-1})\), the generated expression is
\begin{equation}
  N_{\rm auto}(q,p)
  =
  \prod_{e=0}^{|E|-1}
  \prod_{r=0}^{d_e-1}
  \left(q_e\cdot p_{(e+r)\bmod n_{\rm ext}}\right).
  \label{eq:auto-numerator}
\end{equation}
Edges absent from the bound input have \(d_e=0\) and therefore contribute no
factor.  The external index cycle is deterministic and is chosen only to make
the numerator depend on generic external energies; it has no bearing on the
representation builder.  For example, on \texttt{box\_pow3.dot},
\[
  \texttt{--energy-degree-bounds 3:3,4:4 --numerator-expr auto}
\]
is resolved to
\begin{lstlisting}
dot(edges[3], ext[0]) * dot(edges[3], ext[1]) * dot(edges[3], ext[2]) *
dot(edges[4], ext[1]) * dot(edges[4], ext[2]) * dot(edges[4], ext[0]) *
dot(edges[4], ext[1])
\end{lstlisting}
The concrete string is stored in JSON test reports and in Symbolica evaluator
metadata.  For \texttt{evaluate} and \texttt{compile}, \texttt{auto} reads the
bound vector from \texttt{graph.energy\_degree\_bounds} in the input
orientation JSON; for \texttt{build}, \texttt{test}, \texttt{compare}, and
\texttt{test-cff-ltd}, it uses the command-line bound vector.

For repeated-channel interpolation one may additionally set
\begin{lstlisting}
--uniform-numerator-sampling-scale none
--uniform-numerator-sampling-scale beyond-quadratic
--uniform-numerator-sampling-scale all
\end{lstlisting}
on \texttt{build}, \texttt{test}, or \texttt{compare}.  The first value is the
default.  In uniform modes the JSON graph block records
\begin{lstlisting}
"uniform_numerator_sampling_scale": "all",
"uniform_scale_symbol": "M"
\end{lstlisting}
and every energy map is allowed to contain an integer coefficient
\texttt{"m"} multiplying the runtime scale \(M\):
\begin{lstlisting}
{"i": [[3, 1]], "x": [[0, -1]], "m": 2, "c": "0"}
\end{lstlisting}
The expression above means \(\omega_3-E_0+2M\).  Variants also have an integer
\texttt{uniform\_scale\_power}; a value \(r\) multiplies that variant by
\(M^{-r}\).  The arbitrary-precision evaluator receives \(M\) through
\texttt{evaluate --uniform-scale VALUE}.  The \texttt{test} command receives
comparison scales through \texttt{--uniform-scales A,B}; if omitted, it uses
\(M=1.0\) and \(M=2.75\).

The surface cache stores both a compact expression and structured terms:
\begin{lstlisting}
"terms": {
  "ose": [[edge_id, coeff]],
  "external_shift": [[external_id, coeff]],
  "uniform_scale": coeff,
  "constant": "0"
},
"surface_origin": "physical",
"singularity": "physical"
\end{lstlisting}
\texttt{surface\_origin="helper"} is reserved for factors introduced by
interpolation, known-factor reconstruction, finite-pole recursion, or terminal
sectors.  CFF causal surfaces emitted by the original graph or by a
lower-sector CFF component are still classified as
\texttt{surface\_origin="physical"}, because they are genuine causal thresholds
of the denominator set being evaluated.  \texttt{singularity="spurious"} is an
algorithmic marker for a helper factor that is not one of those causal
surfaces, for example a factor containing \(M\), a non-unit OSE coefficient, or
a polynomial reconstruction factor.  In pretty output the class tag adds
\texttt{\_h} for helper origin and \texttt{\_sp} for the spurious marker, while
parentheses still mean numerator-only.  Thus \texttt{e} is a physical
denominator E-surface and \texttt{(e\_h\_sp)} is a numerator-only helper factor.
Pure CFF denominator trees are tested not to contain \texttt{\_sp} surfaces.

After the UV check, the guaranteed support is:
\begin{center}
\begin{tabular}{lll}
\toprule
family & graph class & supported bounds \\
\midrule
\texttt{ltd} & any graph & any UV-convergent bounds; structure unchanged \\
\texttt{hybrid} & no repeated propagators & any UV-convergent bounds; exact LTD collapse \\
\texttt{hybrid} & repeated propagators & any tested UV-convergent bounds; finite-difference samples \\
\texttt{cff} & any graph & all \(d_e\le2\), denominator \(E\)-surfaces only \\
\texttt{cff} & high-power edge caps & channel/contact normal form; denominator \(E\)-surfaces only \\
\bottomrule
\end{tabular}
\end{center}

The last \texttt{cff} row includes one-loop and multiloop sectors with mixed
cubic/quadratic and quartic-or-higher caps, including caps on repeated-signature
edges, provided the loop-energy UV check succeeds.  Repeated channels are
reduced by \eqref{eq:channel-normal-form}; non-repeated high-power contacts are
reduced by the lower-sector recursion
\eqref{eq:lower-sector-normal-form}--\eqref{eq:lower-sector-recursion}.  The
implementation still may become impractical for very large cap vectors because
the number of Lagrange samples and lower-contact sectors grows rapidly, but the
JSON grammar and evaluator semantics do not require a separate fallback branch.

\paragraph{Embedding in gammaLoop cross-section integrands.}
The gammaLoop production path always constructs the three-dimensional
representation through \texttt{generate\_3d\_expr}.  Initial-state cut edges of a
forward-scattering graph are treated as external energy shifts rather than
ordinary CFF/LTD denominator edges.  External tree structures and pure tree
graphs are likewise preserved as residual four-dimensional denominators: when
the reduced source has no loop-energy variables, the energy-bound check is
trivially convergent and the generated object is just the product of those
residual denominators.

Local inspect parity in the cross-section tests compares three structurally
different constructions: CFF with local UV counterterms from 3D expansions, CFF
with local UV counterterms from expanded four-dimensional integrands, and LTD
with local UV counterterms from expanded four-dimensional integrands.  The only
LTD-specific global bridge in this comparison is the loop-measure convention
factor \((-1)^{L-1}\), where \(L\) is the forward-scattering loop count after
removing initial-state cut edges.  CFF expanded-source UV terms instead receive
the difference between the full-graph and reduced-source generalized-CFF sign
exponents; this accounts for duplicate-signature excess lost or gained when a
local UV source is projected out of the original graph.  These signs are
representation-convention factors, not graph-name patches.

The hybrid high-power tests are not pure-CFF tests.  They compare the hybrid
value with a split-mass LTD proxy and verify convergence as the mass split is
reduced.  This is the correct reference for repeated propagators because
ordinary LTD on the unsplit repeated graph would require numerator
derivatives, which the hybrid construction avoids.

\section{Oriented JSON Semantics}

The top-level JSON contains graph metadata, a surface cache, and a list of
orientations.  The central invariant is:
\begin{quote}
one orientation corresponds to one unique EMR edge-energy numerator map.
\end{quote}
All denominator terms sharing that numerator call are stored as
\texttt{variants}.  Each variant has its own rational prefactor, half-edge
list, numerator-surface list, and factorized denominator tree.

The evaluator uses only the JSON plus the DOT graph.  It does not reconstruct a
hidden LTD or CFF object.  A variant contributes
\[
  \texttt{pref}\;
  \prod_{e\in\texttt{half\_edges}}(2\omega_e)^{-1}
  \prod_{s\in\texttt{num\_surfaces}}\eta_s
  \prod_{t\in\texttt{denominator tree}}\eta_t^{-1}
\]
times the numerator sampled at the orientation's \texttt{edge\_q0} map.

\section{Surface Cache}

Each surface cache entry records a class, coefficients of internal on-shell
energies, and external-energy shifts.  The class is printed as
\texttt{e} or \texttt{h}.  If a surface appears only in numerator factors, the
entry carries \texttt{numerator\_only=true} and pretty printing displays
\texttt{(e)} or \texttt{(h)}.  If a numerator-side surface is identical to an
existing denominator surface, the same id is reused.

\section{Orientation Labels and Variants}

Compact labels use \(+\), \(-\), \(0\), and \(x\):
\begin{itemize}[leftmargin=2em]
  \item \(+\): the numerator samples that edge at \(+\omega_e\);
  \item \(-\): the numerator samples that edge at \(-\omega_e\);
  \item \(0\): the numerator samples that edge at zero energy;
  \item \(x\): the edge has a non-trivial affine map, printed by
        \texttt{--show-details}.
\end{itemize}
The physical origin of a contribution, such as CFF, LTD, pinched contacts, or
hybrid basis/interpolation data, is variant metadata rather than part of the
orientation identity.  This is what keeps the number of numerator calls
minimal.

\section{Command-Line Workflow}

The main commands are:
\begin{lstlisting}
python3 generalised_ltd.py validate --dot examples/graphs/box_pow3.dot
python3 generalised_ltd.py build --family hybrid --dot examples/graphs/box_pow3.dot --pretty
python3 generalised_ltd.py build --family cff --dot examples/graphs/box_pow3.dot \
  --energy-degree-bounds 3:3,4:4 --numerator-expr auto
python3 generalised_ltd.py test --dot examples/graphs/box_pow3.dot
python3 generalised_ltd.py test --dot examples/graphs/box_pow3.dot \
  --energy-degree-bounds 0:1,1:1,2:0,3:4 \
  --numerator-expr 'edges[0][0] * edges[1][0] * edges[3][0]**4'
python3 generalised_ltd.py test --dot examples/graphs/box_pow3.dot \
  --energy-degree-bounds 3:3,4:4 --numerator-expr auto
python3 generalised_ltd.py test-cff-ltd --dot examples/graphs/box.dot --energy-degree-bounds 0:2
python3 generalised_ltd.py compile --orientation-json demo.json --dot graph.dot --numerator-expr "1"
python3 generalised_ltd.py evaluate --orientation-json demo.json --dot graph.dot --evaluator-backend symbolica
\end{lstlisting}
The repeated-propagator three-way test compares pure CFF, hybrid, and the
split-mass LTD limiting proxy.  The ordinary LTD formula on an unsplit repeated
graph is not a valid physics comparison because numerator derivatives are not
included there.

\section{Symbolica Evaluators}

The \texttt{compile} command accepts one JSON representation and one fixed
numerator expression.  It reconstructs the JSON expression symbolically,
introducing structured functions such as \texttt{OSE(qx,qy,qz,m)},
\texttt{etaE(i)}, \texttt{etaH(i)}, and tagged edge/loop/external momentum
components.  The resulting Symbolica evaluator is saved both as a shared
library and as eager evaluator state.  The JSON is updated with relative paths
and fingerprints tying the evaluator to the graph, numerator, value type, and
orientation structure.

The \texttt{evaluate} command can use:
\[
\texttt{builtin},\quad
\texttt{symbolica\_compiled},\quad
\texttt{symbolica\_eager},\quad
\texttt{symbolica\_eager\_symjit}.
\]
Profiling performs ten calls of a batch evaluator.  The default batch size is
100, so a profile line reports one thousand repeated samples of the same input
kinematics.

The same saved eager evaluator is used by
\texttt{evaluate --stability}.  This mode converts the generated or supplied
real kinematic input to decimal precision-tracking numbers with 16 significant
digits by default, evaluates with Symbolica's
\texttt{Evaluator.evaluate\_with\_prec}, and reports the number of significant
digits carried by the final decimal result.  The work precision for constants
and arithmetic defaults to 80 decimal digits and can be changed with
\texttt{--stability-work-precision}.  Non-finite outputs are assigned zero
reported digits.  The printed central value is diagnostic only: once the
reported digit count is small, the value may have lost enough significance to
be numerically meaningless.

\chapter{Performance}

\section{Benchmark Setup}

The numbers below were generated on April 26, 2026 using the scripts in
\texttt{examples/scripts}.  Each timing is the per-sample time reported by
\texttt{evaluate --profiling 100}, i.e. ten batch calls of size 100.  The
15-external one-loop benchmark is intentionally omitted here: the LTD build is
small, but the pure CFF one-loop representation has \(2^{15}-2=32766\) sign
sectors and the CFF build is not practical for routine documentation runs.

\section{Repeated Box}

The repeated box script
\texttt{examples/scripts/box\_pow3\_cli\_showcase.sh} compares hybrid and pure
CFF for the six-edge graph with a triple repeated propagator.  The numerator is
\texttt{dot(edges[0], ext[0]) + edges[3][0]}.

\begin{center}
\begin{tabular}{llrrrr}
\toprule
label & backend & orientations & maps & value type & per sample \\
\midrule
hybrid & compiled & 17 & 17 & real & \(0.0955\,\mu s\) \\
hybrid & eager & 17 & 17 & real & \(0.521\,\mu s\) \\
hybrid & eager+symjit & 17 & 17 & real & \(0.104\,\mu s\) \\
hybrid & compiled & 17 & 17 & complex & \(0.271\,\mu s\) \\
hybrid & eager & 17 & 17 & complex & \(0.725\,\mu s\) \\
hybrid & eager+symjit & 17 & 17 & complex & \(0.0989\,\mu s\) \\
CFF & compiled & 62 & 62 & real & \(0.162\,\mu s\) \\
CFF & eager & 62 & 62 & real & \(1.44\,\mu s\) \\
CFF & eager+symjit & 62 & 62 & real & \(0.142\,\mu s\) \\
CFF & compiled & 62 & 62 & complex & \(0.466\,\mu s\) \\
CFF & eager & 62 & 62 & complex & \(2.03\,\mu s\) \\
CFF & eager+symjit & 62 & 62 & complex & \(0.136\,\mu s\) \\
\bottomrule
\end{tabular}
\end{center}

The hybrid representation has fewer numerator maps than CFF, which explains
the faster compiled real evaluator.  Complex compiled evaluation is slower
because the generated function propagates complex arithmetic even though the
input kinematics are real.  Symjit eager evaluation is competitive for this
small expression because the batch is small and the expression has only 17 or
62 maps.

\section{One-Loop 10-Gon}

The script
\texttt{examples/scripts/one\_loop\_10\_external\_runtime\_compare.sh}
compares numerator-free LTD and CFF on a true one-loop 10-gon.

\begin{center}
\begin{tabular}{llrrr}
\toprule
family & backend & orientations & maps & per sample \\
\midrule
LTD & compiled & 10 & 10 & \(0.234\,\mu s\) \\
LTD & eager & 10 & 10 & \(0.852\,\mu s\) \\
LTD & eager+symjit & 10 & 10 & \(0.147\,\mu s\) \\
CFF & compiled & 1022 & 1022 & \(3.96\,\mu s\) \\
CFF & eager & 1022 & 1022 & \(36.6\,\mu s\) \\
CFF & eager+symjit & 1022 & 1022 & \(2.97\,\mu s\) \\
\bottomrule
\end{tabular}
\end{center}

The result is expected: the one-loop CFF orientation count grows like
\(2^n-2\), while LTD has only one cut per internal edge.  The CFF expression is
still fast in absolute terms after compilation, but it is much larger.

\section{Five-Loop Four-External Benchmarks}

The script
\texttt{examples/scripts/five\_loop\_symbolica\_runtime\_compare.sh} compares
three numerator choices on a non-repeated five-loop graph and on the repeated
five-loop graph.  The three numerator choices are \(1\), one edge-external dot
product per internal edge, and the same numerator with three selected edge
terms squared.

\begin{center}
\begin{tabular}{llrrrr}
\toprule
case & family & maps & compiled & eager & eager+symjit \\
\midrule
no repeats, \(N=1\) & LTD & 137 & \(0.326\) & \(1.88\) & \(0.250\) \\
no repeats, \(N=1\) & CFF & 930 & \(1.26\) & \(12.3\) & \(0.892\) \\
no repeats, \(N=1\) & hybrid & 137 & \(0.333\) & \(1.87\) & \(0.246\) \\
no repeats, linear & LTD & 137 & \(0.481\) & \(3.29\) & \(0.379\) \\
no repeats, linear & CFF & 930 & \(1.96\) & \(17.5\) & \(1.35\) \\
no repeats, linear & hybrid & 137 & \(0.484\) & \(3.24\) & \(0.391\) \\
no repeats, squared & LTD & 137 & \(0.493\) & \(3.45\) & \(1.73\) \\
no repeats, squared & CFF & 4692 & \(18.7\) & \(45.9\) & \(13.2\) \\
no repeats, squared & hybrid & 137 & \(0.520\) & \(3.31\) & \(0.390\) \\
\bottomrule
\end{tabular}
\end{center}

\begin{center}
\begin{tabular}{llrrrr}
\toprule
case & family & maps & compiled & eager & eager+symjit \\
\midrule
repeated, \(N=1\) & LTD & 137 & \(0.334\) & \(1.90\) & \(0.272\) \\
repeated, \(N=1\) & CFF & 930 & \(1.26\) & \(12.5\) & \(0.899\) \\
repeated, \(N=1\) & hybrid & 630 & \(0.823\) & \(4.88\) & \(0.579\) \\
repeated, linear & LTD & 137 & \(0.481\) & \(3.44\) & \(0.408\) \\
repeated, linear & CFF & 930 & \(1.88\) & \(17.4\) & \(1.32\) \\
repeated, linear & hybrid & 630 & \(21.4\) & \(37.8\) & \(12.4\) \\
repeated, squared & LTD & 137 & \(0.491\) & \(3.53\) & \(0.387\) \\
repeated, squared & CFF & see text & see text & see text & see text \\
repeated, squared & hybrid & 115 & \(1.84\) & \(9.47\) & \(1.11\) \\
\bottomrule
\end{tabular}
\end{center}

For non-repeated graphs the hybrid family collapses to the LTD structure, so
their orientation and map counts agree.  CFF is slower because it has more
causal sign sectors.  The squared CFF no-repeat case is much larger because the
bounded-degree contact completion adds pinched CFF minors.

For repeated graphs, the LTD rows are included only as runtime references.  The
ordinary unsplit LTD expression does not include numerator derivatives for
raised propagators; the correct comparison is CFF versus hybrid versus the
split-mass LTD limiting test.  The repeated unit numerator shows the expected
middle ground: hybrid has fewer maps than CFF but more than LTD.  The repeated
linear case is slower for hybrid because the confluent interpolation creates
heavier denominator variants per numerator map.  In the squared repeated case,
the pure CFF row now uses the recursive \(E\)-surface-only contact/minor
construction rather than an LTD-contact correction; the timing table should be
regenerated before quoting a stable number for that row.  The bounded hybrid
reduction has 115 maps for this benchmark.

\section{Precision Stability}

The following tables use \texttt{evaluate --stability} with 16 significant
digits in the input and 80 decimal digits of work precision.  They measure the
precision carried by Symbolica's eager precision-tracking arithmetic through the
already compiled structured expression.  These measurements are not identity
tests; they probe numerical conditioning of the particular representation and
kinematic point.  When only one or two digits remain, the displayed value should
be interpreted as a symptom of cancellation rather than as a trustworthy
estimate.

\begin{center}
\begin{tabular}{llrrr}
\toprule
example & family & maps & result digits & tracked value \\
\midrule
repeated box & LTD & 6 & 0 & \(\infty\) \\
repeated box & CFF & 62 & 1 & \(10^{-2}\) \\
repeated box & hybrid & 17 & 1 & \(-2\times 10^{3}\) \\
one-loop 10-gon & LTD & 10 & 2 & \(2.6\times 10^{5}\) \\
one-loop 10-gon & CFF & 1022 & 1 & \(10^{-1}\) \\
\bottomrule
\end{tabular}
\end{center}

For the repeated box, the LTD row is again not a valid repeated-propagator
formula; its singular output is a useful reminder that the unsplit LTD
structure contains zero repeated-channel surfaces.  CFF and hybrid have the
same ordinary double value for this point, but the precision-tracking run shows
that both are badly conditioned at 16 input digits, with the hybrid central
value already unusable.  This is consistent with the hybrid expression being a
more compact but more cancellation-prone confluent combination.

\begin{center}
\begin{tabular}{lrrr}
\toprule
five-loop case & LTD digits & CFF digits & hybrid digits \\
\midrule
no repeats, \(N=1\) & 2 & 2 & 2 \\
no repeats, linear & 2 & 2 & 2 \\
no repeats, squared & 2 & 2 & 2 \\
repeated, \(N=1\) & 0 & 2 & 2 \\
repeated, linear & 0 & 2 & 2 \\
repeated, squared & 0 & 2 & 2 \\
\bottomrule
\end{tabular}
\end{center}

The five-loop tables show that the tested random point is numerically delicate
for all structured sums at ordinary double-like input precision.  The
non-repeated hybrid rows coincide with LTD, as required by construction.  In
the repeated graph, the invalid unsplit LTD expression produces \texttt{NaN}
and therefore zero digits, while CFF and hybrid retain two tracked digits in
all three numerator scenarios.  The result supports using the stability mode as
a representation-level diagnostic: it highlights where algebraically equivalent
JSON structures differ in numerical conditioning even when the ordinary
double-precision evaluator returns matching values.

\begin{thebibliography}{9}
  \bibitem{ltd1906}
  J. J. Aguilera-Verdugo et al.,
  algorithmic multiloop LTD construction, arXiv:1906.06138.

  \bibitem{cff2211}
  Z. Capatti et al.,
  causal/CFF representation construction, arXiv:2211.09653.

  \bibitem{cff2311}
  Z. Capatti et al.,
  further CFF representation and box examples, arXiv:2311.14374.

  \bibitem{cltd}
  \texttt{cLTD} and \texttt{gammaLoop} implementations,
  \url{https://github.com/alphal00p/cLTD} and
  \url{https://github.com/alphal00p/gammaloop}.
\end{thebibliography}

\end{document}
